\documentclass[11pt,oneside,openany]{book}
\usepackage{zlc_frontend_notes}

\begin{document}
\frontmatter
\zlctitlepage{Zou\_lab\_control 主手册}{notebook-first + PyQt GUI + remote FPGA pulse-streamer 系统总览}{架构、session/devices/timing 分层、sequencer 生命周期、真实硬件 runbook 与 N-slot 扫描模型}{2026-06-10}
\tableofcontents
\cleardoublepage
\mainmatter

\chapter{写给新读者}

\section{这套系统是做什么的}
\pyapi{Zou_lab_control} 是一套中性原子实验控制代码库，采用 notebook 优先的工作方式。它把硬件控制、时序、图像分析、绘图和结果对象分层，使得同一份实验逻辑可以在三种后端上运行：

\begin{itemize}
  \item virtual 设备：完全离线，用于教学和回归测试；
  \item 真实 qCMOS 相机；
  \item 远程 FPGA pulse-streamer（脉冲序列器）。
\end{itemize}

整套系统由三块组成，本手册是\tfocus{总览手册}，对应另外两本专题手册：

\begin{description}
  \item[main（本手册）] 系统架构、neutral\_atom 的 session/devices/timing 分层、\pyapi{PulseSequence} 与 \pyapi{PulseTableState} 的区别、sequencer 生命周期、真实硬件 runbook，以及全新的 N-slot 扫描模型。
  \item[frontend] GUI 与绘图：\pyapi{qt_fluent} 控件库、脉冲编辑器三个标签页、按字段扫描点工作流、绘图 API 与 PDF 渲染。见 \filepath{docs/frontend_manual/}。
  \item[fpga] pulse-streamer 硬件：Artix-7 35T 上的 edge-table RTL、1-tick FIFO 预取流水线、2-bank 流式扫描窗口、N-slot 仿射扫描引擎、模拟总线 DAC 引擎、JTAG-to-AXI 上传流程与资源预算。见 \filepath{docs/fpga_manual/}。
  \item[device] 设备与实验：设备配置/加载、相机采集、相机读出原理（sitemap/thresholds/detect）、标定与结果对象、完整实验流程。见 \filepath{docs/device_manual/}。
\end{description}

\begin{notebox}[阅读约定]
本手册用中性教学口吻。每个组件归属哪一层、状态住在哪里、调用后应得到什么结果，都会写清楚。维护者用的架构不变式和反模式放在 \filepath{docs/MAINTAINER_NOTES.md}，不放进本手册。
\end{notebox}

\section{三台角色的物理图景}
真实实验是双电脑结构：一台控制/qCMOS 电脑，一台 FPGA/Vivado 电脑，两者用 RPyC 连接。

\begin{figure}[htbp]
  \centering
  \begin{tikzpicture}[
      node distance=6mm,
      box/.style={draw=StarkGray, rounded corners=2pt, fill=StarkPale, align=center, inner sep=4pt, font=\footnotesize, minimum height=10mm, text width=30mm},
      hw/.style={box, fill=StarkTint},
      arr/.style={-{Latex[length=2mm]}, StarkTeal, thick}]
    \node[box] (nb) {控制/qCMOS 电脑\\Jupyter / Pulse GUI};
    \node[box, right=12mm of nb] (remote) {RemoteSequencer\\(RPyC client)};
    \node[box, right=12mm of remote] (server) {SequencerService\\FPGA 电脑 server};
    \node[hw, below=10mm of server] (vivado) {Vivado hw\_axi 会话\\persistent Tcl (JTAG-to-AXI)};
    \node[hw, below=10mm of remote] (rtl) {zlc\_pulse\_streamer\_top\\edge-table 引擎};
    \node[hw, below=10mm of nb] (out) {62 路 TTL\\+ 4x10b DAC\\+ emCCD 触发};
    \draw[arr] (nb) -- (remote);
    \draw[arr] (remote) -- (server);
    \draw[arr] (server) -- (vivado);
    \draw[arr] (vivado) -- (rtl);
    \draw[arr] (rtl) -- (out);
  \end{tikzpicture}
  \caption{真实硬件路径。GUI/notebook 只生成脉冲状态；编译、上传和 fire 由 FPGA 电脑负责。}
  \label{fig:main-arch}
\end{figure}

GUI 只是脉冲前端。它可以只显示 \pyapi{trap/cooling/probe/emCCD} 四路，也可以显示全部 XDC 推断出来的 channel；上传时永远按 server 的完整硬件 channel order 编译，未显示或未配置的 channel 补成 off。

\chapter{原理：分层与状态归属}

\section{neutral\_atom 的分层}
\pyapi{Zou_lab_control.neutral_atom} 围绕清晰的边界组织：

\begin{description}
  \item[\pyapi{devices/}] 硬件适配器与设备契约。设备只负责硬件动作。包含 \filepath{axi_session.py}（JTAG-to-AXI 上传 \pyapi{VivadoAxiStreamerSession}）、\filepath{fpga_pulse_streamer.py}（XDC 推断 + 程序校验）、\filepath{sequencer.py}、\filepath{sequencer_server.py}。
  \item[\pyapi{timing/}] 时序的真相来源。\filepath{sequence.py} 里的 \pyapi{PulseSequence}，\filepath{pulse_table.py} 里的 \pyapi{PulseTableState}（GUI 编译模型），\filepath{verilog.py} 的边沿表生成。
  \item[\pyapi{operations/}] 纯图像/标定/检测算法，可离线运行。
  \item[\pyapi{subsystems/}] 实验级工作流，例如 \pyapi{exp.readout}、\pyapi{exp.timing}。
  \item[\pyapi{views/}] 接到 frontend 的绘图适配器。
\end{description}

\begin{notebox}[一个关键不变式]
相机 capture 只显示\tfocus{原始}图像。site 圆圈、阈值这些标定叠加属于 readout 结果，不属于相机设备。如果 capture 上出现了 site 圆圈，那是把标定状态错放到了相机层。
\end{notebox}

\section{连接实验}
\pyapi{na.connect} 返回一个 \pyapi{NeutralAtomSession}。三种典型连法：

\begin{codeblock}[Python]
import Zou_lab_control.neutral_atom as na

# 完全离线：虚拟相机 + 虚拟 sequencer
exp = na.connect("virtual")

# 真实硬件：控制电脑连 FPGA 电脑上的 server
exp = na.connect(
    "remote_template",
    sequencer={"host": "192.168.0.20", "port": 18861},
    open_devices=True,
)
\end{codeblock}

session 暴露子系统：\pyapi{exp.camera}（采图）、\pyapi{exp.timing}（脉冲/触发）、\pyapi{exp.readout}（标定与占据检测）、\pyapi{exp.devices}（底层设备，例如 \pyapi{exp.devices.sequencer}）。

\chapter{PulseSequence 与 PulseTableState}

\section{两个模型的分工}
时序有两个层次的对象，不要混淆：

\begin{description}
  \item[\pyapi{PulseSequence}] 底层、与时间相关的真相来源。它知道每条 channel 在每个时刻的电平，能编译成边沿表（绝对 tick + 完整输出 mask）。notebook、PyQt、远程 FPGA 都共享这同一个真相来源。
  \item[\pyapi{PulseTableState}] GUI 的\tfocus{编译模型}。它保留用户能理解的概念：一排\tfocus{周期 (period)}，每个 period 有一个持续时间和一整套数字状态向量；再加上 channel 名称/标签、可见性、延时、模拟总线计划，以及扫描 slot。它本身不拥有硬件，只 \pyapi{compile()} 成 \pyapi{PulseSequence}/runtime program。
\end{description}

一句话：\pyapi{PulseTableState} 是给人编辑的表格，\pyapi{PulseSequence} 是给时钟执行的边沿。

\section{从状态到边沿表}
\begin{codeblock}[Python]
state = na.PulseTableState.load("pulses/camera_imaging_address_switch.json")

program = state.compile(
    clock_hz=exp.devices.sequencer.clock_hz,        # 默认 50 MHz
    trigger_channels=exp.devices.sequencer.trigger_channels,
    repeat_forever=False,
)
print(program.ticks)   # 例如 [0, 100000, 105000, 106000, 1105000, 1106000]
print(program.masks)   # 每个 tick 对应的完整输出 mask
\end{codeblock}

一行边沿的含义是：\tfocus{在这个绝对 FPGA tick，把所有输出切换成这个 mask}。bit 0 对应 \pyapi{channels[0]}，bit 1 对应 \pyapi{channels[1]}，以此类推。多条 channel 在同一物理时刻变化，会合并成一条带完整 mask 的边沿。

\begin{codeblock}[Camera 预设在 50 MHz 的编译结果]
ticks: 0, 100000, 105000, 106000, 1105000, 1106000
masks: 513, 512, 2568, 520, 512, 0
\end{codeblock}

\pyapi{513 = ch00 + ch09}（cooling + trap），\pyapi{2568 = ch03 + ch09 + ch11}（probe + trap + emCCD）。相机外触发 bit 是 \pyapi{ch11/emCCD}。

\chapter{Sequencer 生命周期}

\section{prepare / fire / wait\_done / safe\_state}
控制电脑通过 \pyapi{RemoteSequencer} 调用四个动作。它们经 RPyC 到达 FPGA 电脑上的 \pyapi{SequencerService}，再到 \pyapi{VivadoAxiStreamerSession}：

\begin{figure}[htbp]
  \centering
  \begin{tikzpicture}[
      stage/.style={draw=StarkTeal, rounded corners=2pt, fill=StarkCodeBack, align=center, font=\footnotesize, minimum height=14mm, text width=26mm, inner sep=3pt},
      arr/.style={-{Latex[length=2mm]}, StarkBlue, thick}]
    \node[stage] (prep) {\textbf{prepare}\\SAFE + 打包上传\\BRAM image\\(edges/scan/bus)\\发 LOAD};
    \node[stage, right=8mm of prep] (fire) {\textbf{fire}\\发 FIRE\\(COMMAND 上升沿)};
    \node[stage, right=8mm of fire] (wait) {\textbf{wait\_done}\\轮询 STATUS\\流式回填扫描 bank};
    \node[stage, right=8mm of wait] (safe) {\textbf{safe\_state}\\复位输出\\回到安全态};
    \draw[arr] (prep) -- (fire);
    \draw[arr] (fire) -- (wait);
    \draw[arr] (wait) -- (safe);
  \end{tikzpicture}
  \caption{一次 shot 的生命周期。prepare 只是把表写好并校验，不启动；只有 fire 的同步上升沿才是 start 事件。}
  \label{fig:lifecycle}
\end{figure}

\begin{description}
  \item[\pyapi{prepare(program)}] 先发 SAFE，把编译好的程序打包成 BRAM image（边沿表、重复元数据、扫描点、总线段），经 JTAG-to-AXI 写进 FPGA，arm 扫描 bank，然后发 LOAD（COMMAND 上升沿，等 \pyapi{STATUS_LOADED}）。\tfocus{prepare 不启动脉冲}，它只是让 FPGA 持有一张已校验的表。
  \item[\pyapi{fire()}] 发 FIRE 命令（COMMAND 字先清 0 再置位，制造干净的上升沿）。FPGA 只把同步后的上升沿当作 start。这之后，微秒级时序由 FPGA 时钟拥有；Python、RPyC、Windows 调度、Vivado、JTAG 延迟都不再决定 shot 内的边沿。
  \item[\pyapi{wait_done()}] 轮询 \pyapi{STATUS}，等待有限序列跑完；对超出常驻窗口的流式扫描，它会跟随 \pyapi{CURSOR} 把空出来的扫描 bank 回填下一段。
  \item[\pyapi{safe_state()}] 发 SAFE，把输出拉回安全态（\pyapi{Stop Pulse} 走这条路）。
\end{description}

\section{Vivado 持久会话与重复 prepare 缓存}
\pyapi{jtag-axi} 后端只启动一次 Vivado Tcl 进程，打开 hardware target，加载一次 \pyapi{.ltx}（让 Vivado 在已烧录的比特流里找到 \pyapi{jtag_axi} 核），然后整个 server 生命周期内复用这个 \pyapi{hw_axi} 会话来做 prepare/fire/wait\_done/safe\_state。慢的 hardware-target 建立发生在 server 启动时，而不是第一次 \pyapi{On Pulse}。

\pyapi{SequencerService.prepare} 还做了一层缓存：如果新程序的 \pyapi{sequence_id}（程序内容的哈希）和已准备的完全相同，就跳过整段上传，直接复用已经载入的程序。是否启用由 \pyapi{ZLC_SEQUENCER_CACHE_PREPARED} 控制。所以一次扫描里若某点和上一点编译结果一致，第二次 \pyapi{On Pulse} 只是几次控制写。

\section{API 等价写法}
\begin{codeblock}[Python]
state = na.PulseTableState.load("pulses/camera_imaging_address_switch.json")
program = state.compile(
    clock_hz=exp.devices.sequencer.clock_hz,
    trigger_channels=exp.devices.sequencer.trigger_channels,
    repeat_forever=False,
)
exp.devices.sequencer.prepare(program)
exp.devices.sequencer.fire()
exp.devices.sequencer.wait_done()
\end{codeblock}

\pyapi{repeat_forever=True} 适合示波器找线，输出会一直循环。真实相机采集有限帧时，应让 readout/camera helper 先 arm 相机，再生成刚好需要帧数的有限触发序列，不要让 qCMOS 等一个没有帧数边界的自由循环脉冲。

\chapter{真实硬件 Runbook}

\section{安装}
在两台电脑上都装 Python 包；Vivado 只装在 FPGA 电脑。

\begin{codeblock}[PowerShell]
cd D:\ZLC
install_requirements.bat
\end{codeblock}

\filepath{install_requirements.bat} 把当前解释器记到 \filepath{.zlc_python_path}；之后 \filepath{pulse_gui.bat}、\filepath{fpga\run_server.bat} 优先用同一个解释器。若 Vivado 不在默认路径：

\begin{codeblock}[PowerShell]
$env:ZLC_PS_VIVADO_BIN = "C:\Xilinx\Vivado\2019.1\bin\vivado.bat"
\end{codeblock}

\section{在 FPGA 电脑 build 和 program}
\begin{codeblock}[PowerShell]
cd D:\ZLC
.\fpga\build_and_program.bat --check    # 不连板的 HDL 综合 + 容量自查
.\fpga\build_and_program.bat            # build + program (create_project.tcl -> program_fpga.tcl)
\end{codeblock}

默认工程在 \filepath{fpga\build\ps}（短名 \pyapi{ps} 让 Vivado 深层 run/.Xil 临时路径不超过 Windows MAX\_PATH，构建仍留在仓库内），bit/ltx（\filepath{zlc_pulse_streamer_top.bit/.ltx}）在其 \filepath{ps.runs\textbackslash impl_1} 目录下。默认 XDC 是 address-switch 板的引脚图（\filepath{fpga\textbackslash board_config\textbackslash board.xdc}，见该目录 README；用 \pyapi{$env:ZLC_PS_XDC} 覆盖）。相机成像预设的关键物理输出：

\begin{codeblock}[Camera-imaging subset]
ch09 -> trap    -> M17
ch00 -> cooling -> F15
ch03 -> probe   -> N15
ch11 -> emCCD   -> M13   (qCMOS 外触发)
\end{codeblock}

默认时钟 50 MHz，最小步进 20 ns。如果示波器看到脉冲长度是设定值的两倍，先检查 control/server/GUI 是否还在沿用旧的 100 MHz 假设。

\section{启动 server}
\begin{codeblock}[PowerShell]
.\fpga\run_server.bat --check-config    # 打印 project/bit/ltx/XDC/channel/trigger/clock/容量
.\fpga\run_server.bat                    # host 0.0.0.0, port 18861
\end{codeblock}

\section{控制电脑 notebook}
\begin{codeblock}[Python]
exp = na.connect("remote_template",
                 sequencer={"host": "192.168.0.20", "port": 18861},
                 open_devices=True)

exp.timing.configure_imaging(exposure=2e-3, load=True)
exp.timing.preflight().raise_if_failed()

capture   = exp.camera.capture(frames=1)
sitemap   = exp.readout.sitemap(frames=20, grid_shape=(5, 7))
threshold = exp.readout.thresholds(frames=120, site=0)
shot      = exp.readout.detect()
\end{codeblock}

\chapter{N-slot 扫描模型}

\section{为什么是“命名 slot”}
这是当前唯一的扫描概念。它取代了旧的 \pyapi{x}/\pyapi{y} 二变量仿射扫描——\tfocus{已经没有 x/y 的说法了}。

任何\tfocus{逐字段}的值都可以\tfocus{绑定到一个命名 slot}：

\begin{itemize}
  \item 一个 period 的持续时间（\pyapi{kind="duration"}）；
  \item 一个 period 内某条模拟总线的 DAC 值（\pyapi{kind="dac"}）。
\end{itemize}

只有这\tfocus{两类}字段可以扫描：\pyapi{SCAN_SLOT_KINDS = ("duration", "dac")}。\tfocus{channel 延时不是扫描量}——它是每条 channel 的一个\tfocus{固定}值（见下文“固定延时如何与扫描共存”），\pyapi{state.bind_field("delay", ...)} 会直接抛 \pyapi{ValueError("delay is a fixed per-channel value and cannot be scanned")}。

slot 按绑定顺序命名为 \pyapi{s0, s1, ...}。在 GUI 里点字段旁边的扫描圆点，或在 API 里调用 \pyapi{state.bind_field(kind, target)} 完成绑定。绑定后该字段的值在表达式里变成 \pyapi{s0} 这样的符号。

\section{scan\_table：N\_points x N\_slots}
扫描数据是一张 \pyapi{scan_table}：行 = 一个扫描点，列 \pyapi{j} = slot \pyapi{s\{j\}} 在该 slot 显示单位下的取值（时间 slot 用 ns，\pyapi{dac} slot 用整数 DAC code）。它可以从 \filepath{.npy}/\filepath{.csv}/\filepath{.txt} 载入，或在 GUI 的 Scan 标签页用 Python 现场构造。

\begin{figure}[htbp]
  \centering
  \begin{tikzpicture}[font=\scriptsize]
    \matrix (m) [matrix of nodes, nodes={draw=StarkPale, minimum width=12mm, minimum height=6mm, anchor=center}, column sep=-\pgflinewidth, row sep=-\pgflinewidth,
      row 1/.style={nodes={fill=StarkTint, font=\bfseries}}] {
      |[fill=StarkPale]| point & s0 (ns) & s1 (ns) \\
      0 & 1000 & 0 \\
      1 & 2000 & 500 \\
      2 & 4000 & 1000 \\
    };
    \node[right=4mm of m, text width=42mm, align=left, font=\footnotesize] {主机编译成\\一张仿射边沿模板\\+ 一张流式扫描点表；\\硬件无缝迭代每个点。};
  \end{tikzpicture}
  \caption{一张两 slot、三点的 scan\_table。}
  \label{fig:scan-table}
\end{figure}

\section{主机如何编译扫描}
绑定过的 duration 是\tfocus{时间 slot}，它参与仿射 tick 公式。主机把每行边沿编译成一个 base tick 加 \pyapi{NUM_SLOTS} 个定点系数，FPGA 在运行时计算：

\begin{equation}
\text{effective\_tick} = \text{base\_tick} + \left(\sum_j \text{coeff}_j \cdot \text{slot}_j\right) \gg \text{COEFF\_FRAC\_BITS}
\label{eq:affine}
\end{equation}

其中 \pyapi{COEFF_FRAC_BITS=8}（定点小数位）。这样一个很大的一维或多维扫描\tfocus{不会}让边沿表变长——边沿数限制的是模板复杂度，不是扫描点数。

时间表达式只能是 slot 的\tfocus{仿射}函数（数字、\pyapi{s0..}、加减乘除、括号）。编译器会拒绝非仿射的扫描时序，也会拒绝那些在不同扫描点上边沿顺序会交换的模板（此时应拆成多个 chunk，或每个点单独 prepare）。

\section{端到端的可运行示例}
下面扫描 period 0 的持续时间（slot \pyapi{s0}）和 period 2 上 \pyapi{da\_dipole} 总线的 DAC 值（slot \pyapi{s1}），同时给 \pyapi{probe} channel 设一个\tfocus{固定}延时（500\,ns，不是扫描量）：

\begin{codeblock}[Python]
import numpy as np
import Zou_lab_control.neutral_atom as na

state = na.PulseTableState.load("pulses/camera_imaging_address_switch.json")

# probe 的延时是一个 FIXED 的每通道值 —— 不进扫描表
state.delays["probe"] = 500.0                 # 固定 500 ns 输出延时

# 绑定两个可扫字段 -> s0, s1（bind_field 是幂等的，返回 slot 序号）
s0 = state.bind_field("duration", "0")        # period 0 的持续时间
s1 = state.bind_field("dac", "da_dipole@2")   # period 2 上 da_dipole 的 DAC 码

# 构造 5 个扫描点的 (N_points x 2) 表：第一列是 ns，第二列是带符号 DAC 值（0 = 0 V）
points = np.linspace(1_000, 5_000, 5)
codes = np.linspace(-512, 511, 5).astype(int)
scan_table = np.column_stack([points, codes])
state.set_scan_table(scan_table)              # 也可 state.load_scan_table("scan.npy")

# 编译成扫描 program 并上传
program = state.compile_scan(clock_hz=exp.devices.sequencer.clock_hz)
exp.devices.sequencer.prepare(program)
exp.devices.sequencer.fire()
exp.devices.sequencer.wait_done()
\end{codeblock}

编译产物 \pyapi{RuntimeSequenceProgram} 在 schema 里携带 \pyapi{slot_count}、\pyapi{slot_kinds}、\pyapi{tick_slot_coeffs}、\pyapi{loop_end_slot_coeffs}、\pyapi{scan_points} 和 \pyapi{scan_coeff_frac_bits}，再加上常规的 ticks/masks/bus\_segments。FPGA 一次上传后无缝迭代所有扫描点，每点对外触发一次相机。

\section{扫描 DAC 值（硬件无缝）}
模拟总线（DAC）的值也能逐点无缝扫描，和数字时序扫描共用同一张扫描点表，不需要每点重新 prepare。绑定 \pyapi{kind="dac"}，\pyapi{scan_table} 的对应列填\tfocus{原始整数 DAC code}（不是电压、不是 ns）：

\begin{codeblock}[Python]
import numpy as np
import Zou_lab_control.neutral_atom as na

state = na.PulseTableState.load("pulses/camera_imaging_address_switch.json")

# 把 da_dipole 在 period 2 的 DAC 电平绑成一个扫描 slot。
sd = state.bind_field("dac", "da_dipole@2")        # target = "<bus>@<period_index>"

# 11 个点，带符号 10-bit DAC 值（-512..+511，0 = 0 V）。
state.set_scan_table(np.linspace(0, 1000, 11).astype(int).reshape(-1, 1))

program = state.compile_scan(clock_hz=exp.devices.sequencer.clock_hz)
exp.devices.sequencer.prepare(program)
exp.devices.sequencer.fire()                       # 逐点切 DAC，每点触发一次相机
\end{codeblock}

底层上，被扫的 DAC 段在边沿模板里不占数字边沿，而是变成一个带 \pyapi{value_select} 的总线段：FPGA 在每个扫描点从对应 slot 读 DAC code 写到总线输出。DAC slot 的列在扫描点表里存原始 code（不做 ns$\to$tick 换算），其仿射系数恒为 0，所以它从不进入 effective\_tick 公式。机制细节见 FPGA 手册“无缝扫描 DAC 值”一节。

\begin{notebox}[组合规则]
\tfocus{延时不是扫描量}——它是每条 channel 的一个\tfocus{固定}值，不进扫描表；它会作为一条加在输出上的\tfocus{逐通道延时线}\tfocus{与任意扫描共存}（duration 扫描、DAC 扫描都行）。可以扫的只有 \tfocus{DAC 值}与\tfocus{持续时间}，两者\tfocus{可以同时扫}（不再互斥）。被扫电平要用 \pyapi{edge} 模式保持，不要用 \pyapi{ramp} 端点去扫。延时（正/负/零）大小\tfocus{有界}：由一条深度 \pyapi{DELAY_DEPTH}（默认 2048 节拍 $\approx\pm40.96\,\mu$s，覆盖 $\pm15\,\mu$s）的延时线实现，对 TTL 和 DAC 都一样——详见下一节。
\end{notebox}

\section{固定延时如何与扫描共存（有界深度的延时线）}
channel 延时\tfocus{不可扫描}：\pyapi{SCAN_SLOT_KINDS = ("duration", "dac")}，\pyapi{state.bind_field("delay", ...)} 会抛 \pyapi{ValueError("delay is a fixed per-channel value and cannot be scanned")}。延时是每条 channel 的一个\tfocus{固定}值，通过 \pyapi{state.delays[channel] = ...}（配 \pyapi{state.delay_units}）设定，编译时作为常数携带（\pyapi{channel_delays}）。

它在硬件上被实现成一条加在引擎\tfocus{输出}上的\tfocus{延时线}：\pyapi{output_delayed[t] = output_undelayed[t-d]}，\pyapi{t < d} 之前恒为 0。它有三条关键性质：

\begin{itemize}
  \item \tfocus{大小有界（覆盖常用区间）}。延时 $d$ 可以是正、负、零，且\tfocus{与帧长无关}（可大于一帧、可小于一个 tick）；唯一的上限是延时线的深度 \pyapi{DELAY_DEPTH}（默认 2048 节拍 $\approx\pm40.96\,\mu$s，覆盖 $\pm15\,\mu$s 的常用区间并留有余量）。负延时用一个全局平移 $G$ 折算成全部为非负，所以真正的约束是一组延时的\tfocus{跨度} $\le$ \pyapi{DELAY_DEPTH}；超出时主机在 \pyapi{validate}/打包阶段（上硬件\tfocus{之前}）抛出指明 \pyapi{DELAY_DEPTH} 的清楚错误，GUI 也会把单个输入夹回上限。
  \item \tfocus{物理延时、第一帧真实}。这是真实的物理延时，不是 GUI 预览里那种循环 \%$T$ 旋转：fire 之后那条 channel \tfocus{安静}到 $t=d$ 才开始，\tfocus{第一帧是真实的}（没有"绕回来的尾巴"幻影）。负延时把整个序列重新平移一个全局量 $G$。
  \item \tfocus{绝不打扰别的 channel、与扫描均匀组合}。延一条 channel 不改变任何其它 channel 的时序，也不改变周期长度；延时与 duration 扫描、DAC 扫描、内层 repeat bracket \tfocus{均匀组合}——"延时后的序列"恒等于"把那条 channel 延时 $d$ 的无延时序列"。
\end{itemize}

\tfocus{DAC 总线有完全相同的能力}：一个 DAC 值（固定、被扫、或斜坡 ramp）都能用同样的方式延时（同一深度上限），$t<d$ 之前保持安全值，延时\tfocus{超过一帧}也支持；一条 DAC 总线的 10 个 bit \tfocus{共享一个}延时。机制细节（环形缓冲、启动门、有界深度）见 FPGA 手册"逐通道输出延时线"一节。

\section{扫描 + 采集：一次完整循环}
硬件无缝扫描时，FPGA 在一次 \pyapi{fire()} 里跑完所有扫描点，每点对外触发一次相机。所以典型用法是：arm 相机收 \pyapi{N_points} 帧，prepare/fire 扫描 program，再按点 reshape 结果。

\begin{codeblock}[Python]
import numpy as np
import Zou_lab_control.neutral_atom as na

exp = na.connect("remote_template", sequencer={"host": "192.168.0.20", "port": 18861}, open_devices=True)
state = na.PulseTableState.load("pulses/camera_imaging_address_switch.json")

s0 = state.bind_field("duration", "1")               # 扫成像段时长
points = np.linspace(1_000, 5_000, 21)               # 21 个点, ns
state.set_scan_table(points.reshape(-1, 1))

program = state.compile_scan(clock_hz=exp.devices.sequencer.clock_hz)
exp.camera.arm(frames=len(points))                   # 先 arm 刚好这么多帧
exp.devices.sequencer.prepare(program)
exp.devices.sequencer.fire()                         # 无缝跑完 21 个点
images = exp.camera.read()                           # (21, H, W)
\end{codeblock}

要点：用\tfocus{有限}扫描 program（不要 \pyapi{repeat_forever}），让相机的帧数边界和扫描点数对齐；不要让 qCMOS 等一个没有帧数边界的自由循环脉冲。

\section{解绑与单点求值}
\begin{codeblock}[Python]
state.unbind_slot(s1)                 # 删除 slot；后面的 slot 自动顺移 (s2 -> s1, ...)

# 不做硬件扫描、只想用一组常数跑一次时：
one_shot = state.with_slots_resolved({"s0": 2_000})   # 把 s0 替换成 2000 ns 的常数副本
\end{codeblock}

\chapter{分层与状态归属（深入）}

本章把整个 \filepath{Zou_lab_control.neutral_atom} 控制框架拆开，逐层讲清楚\tfocus{每一层负责什么、每一份状态归谁所有、谁有权改它、生命周期是怎样的}。这套系统看上去很大，但它的设计哲学其实只有一句话：\tfocus{时间真相、硬件动作、纯算法、实验编排各居其位，互不越界}。读完本章，一个第一次接触这个仓库的人应该能够：理解每一层的职责边界、用 \pyapi{na.connect} 拿到一个会话对象、知道这个对象里挂着哪些状态、在三种后端（虚拟/本地/远程）之间切换、并亲手跑通一次最小实验。

\section{为什么要分层：一句话的设计哲学}

控制一个中性原子实验，本质上要同时管好四类完全不同的东西：

\begin{itemize}
\item \tfocus{时间}：什么通道在什么时刻翻转、脉冲多长、扫描参数怎么平移每条边沿。这是确定性的、可编译的"真相"。
\item \tfocus{硬件}：相机怎么曝光、FPGA 怎么收脉冲表、JTAG-to-AXI 怎么写 BRAM 和控制寄存器。这是有副作用、会失败、需要连接的"动作"。
\item \tfocus{算法}：从一张原始图里找出原子、做校准拟合、判定每个格点亮不亮。这是纯函数式的、可离线复现的"计算"。
\item \tfocus{编排}：把上面三者串成"准备序列 → 开火 → 读出 → 分析 → 存历史"的实验循环。这是"流程"。
\end{itemize}

如果把这四类混在一起（比如让相机驱动顺手画出格点圆圈、让设备类里塞拟合算法），系统很快就会变得无法测试、无法离线复现、无法换后端。因此本框架把它们\tfocus{物理隔离}成不同的子包，并用几条铁律（见本章最后一节"不变量"）锁死边界。

\section{六个子包：职责与状态归属}

\filepath{Zou_lab_control/neutral_atom/} 下分成六个目录。下面逐个说明它"做什么"和"拥有什么状态"。

\subsection{devices/ —— 硬件适配器与契约}

\begin{description}
\item[做什么] 把每一台真实仪器包装成一个统一接口的 Python 对象，\tfocus{只做硬件动作}：连接、配置、触发、读回字节。它不懂算法，也不懂实验流程。
\item[关键文件] \filepath{axi_session.py}（\pyapi{VivadoAxiStreamerSession}：JTAG-to-AXI 打包上传 + CTRL mailbox），\filepath{fpga_pulse_streamer.py}（XDC 推断 + \pyapi{validate_pulse_streamer_program}），\filepath{sequencer.py}（\pyapi{RuntimeSequenceProgram}、\pyapi{SequencerService}、\pyapi{RemoteSequencer}），\filepath{sequencer_server.py}（\pyapi{run_server}、\pyapi{CommandSequencerBackend}），\filepath{qcmos.py}（qCMOS 相机适配器），\filepath{virtual.py}（离线假设备），\filepath{registry.py}（设备注册/查找），\filepath{base.py}（设备基类与契约）。
\item[拥有的状态] 只有\tfocus{与硬件连接相关}的状态：打开的句柄、hw\_axi 会话、当前已写入硬件的脉冲程序、相机的曝光参数。它\tfocus{不}持有图像分析结果，也\tfocus{不}持有校准。
\end{description}

这里有一条贯穿全书的契约：\tfocus{\pyapi{TrapArrayDevice} 暴露 \pyapi{n_sites}（有多少个格点），但绝不暴露每个格点在相机上的像素中心}。像素中心是几何/光学事实，属于校准（\pyapi{TrapCalibration}），不属于设备。设备只知道"我这套光阱有 N 个位点"，至于这 N 个点落在相机的哪些像素上，要去问校准。

\subsection{timing/ —— 时间真相}

\begin{description}
\item[做什么] 定义脉冲序列的"真相"，并把它编译成 FPGA 能吃的边沿表。这是整个系统里唯一允许定义"时间是什么"的地方。
\item[关键文件] \filepath{sequence.py} 里的 \pyapi{PulseSequence}（程序化构造序列的核心数据结构）；\filepath{pulse_table.py} 里的 \pyapi{PulseTableState}（GUI 用的编译模型，承载周期/通道延迟/DAC 总线段/扫描槽 \pyapi{scan_slots} 与扫描表 \pyapi{scan_table}）；\filepath{verilog.py}（生成微型边沿表）。
\item[拥有的状态] 当前序列的全部时间定义：周期、每通道延迟、总线段、扫描槽绑定与扫描点表。它\tfocus{不}碰硬件，编译产物（\pyapi{RuntimeSequenceProgram}）交给 devices 层去上传。
\end{description}

时间层是\tfocus{纯数据 + 纯编译}的。你可以在完全没有硬件、没有 Vivado 的机器上构造一个 \pyapi{PulseSequence}、调用编译、检查产物，这正是单元测试和离线开发的基础。

\subsection{operations/ —— 纯算法（离线）}

\begin{description}
\item[做什么] 收纳所有\tfocus{纯函数式}的图像处理、校准拟合、原子检测算法。给定输入数组，返回结果，\tfocus{没有任何副作用}，不连接任何硬件。
\item[拥有的状态] 几乎没有。它是算法的家，不是数据的家。输入从外面传进来，结果返回出去。
\end{description}

把算法关在这里有两个巨大好处：一是\tfocus{可离线复现}——拿一张存档的原始图，离线就能重跑检测，得到逐位元一致的结果；二是\tfocus{可单测}——算法测试不需要相机、不需要网络。

\subsection{subsystems/ —— 实验级编排}

\begin{description}
\item[做什么] 把 devices 的硬件动作和 operations 的算法\tfocus{编排}成实验级能力。这一层是会话对象上 \pyapi{exp.readout}、\pyapi{exp.timing} 背后的实现。
\item[关键成员] \pyapi{ReadoutSubsystem}（读出子系统：触发相机、拿原始图、调 operations 做检测、产出每格点占据判定），\pyapi{TimingSubsystem}（时序子系统：在会话级别管理当前序列的准备与开火）。
\item[拥有的状态] 编排所需的中间状态与对底层设备/算法的引用，而非原始硬件句柄本身。
\end{description}

\subsection{views/ —— 给前端的绘图适配器}

\begin{description}
\item[做什么] 把内部数据结构转换成前端（GUI/笔记本）能直接画的形式。它是"展示适配器"，不产生新数据。
\item[拥有的状态] 无持久状态；它是从模型到可视化的纯转换层。
\end{description}

\subsection{core/ —— 校准与结果模型}

\begin{description}
\item[做什么] 定义跨实验复用的核心模型与分析入口。
\item[关键文件] \filepath{calibration.py} 里的 \pyapi{TrapCalibration}（把格点编号映射到相机像素几何，即前面说的"像素中心"的真正归属地），\filepath{analysis.py}（分析流程），\filepath{results.py}（结果数据结构）。
\item[拥有的状态] 校准对象本身是\tfocus{一份可被会话持有、可保存/加载}的状态。它是连接"设备说有 N 个格点"和"图像上这 N 个点在哪"的桥梁。
\end{description}

\begin{notebox}[一句话记住分层]
\tfocus{timing} 说"时间是什么"，\tfocus{devices} 说"怎么对硬件动手"，\tfocus{operations} 说"怎么从数据算出答案"，\tfocus{subsystems} 把它们编排成实验能力，\tfocus{core} 提供校准与结果模型，\tfocus{views} 负责画出来。没有任何一层越界去做别人的事。
\end{notebox}

\section{na.connect 与会话对象}

整个框架对用户暴露的入口非常窄：你调用 \pyapi{na.connect(config)}，拿到一个 \pyapi{NeutralAtomSession}。\tfocus{这个会话对象就是你与整个实验交互的唯一句柄}——所有状态都挂在它身上，所有操作都从它出发。

\begin{codeblock}[Python]
import Zou_lab_control.neutral_atom as na

# config 选择后端与硬件参数；下一节详述三种后端
session = na.connect(config)
\end{codeblock}

\subsection{会话拥有哪些状态}

\pyapi{NeutralAtomSession} 是这套系统的\tfocus{状态归属中心}。它持有以下几份状态，每一份都有明确的所有者语义：

\begin{description}
\item[devices] 一个 \pyapi{DeviceSet}：本次会话连接的全部设备（相机、序列器等）的集合。设备的硬件句柄归它管。
\item[sequence] 当前的 \pyapi{PulseSequence}：本会话此刻的时间真相。你修改序列、编译序列，改的都是这一份。
\item[\_calibration] 一个 \pyapi{TrapCalibration} 或 \pyapi{None}：当前校准。下划线前缀表示它是受控状态——通过会话的接口设置/使用，而不是随便外部赋值。没校准时为 \pyapi{None}。
\item[history] 一个列表：本会话历次实验产出的结果，按时间追加。它是你这次实验跑了什么的"账本"。
\item[readout] 一个 \pyapi{ReadoutSubsystem}：读出能力的实现，即 \pyapi{exp.readout} 背后的对象。
\item[timing] 一个 \pyapi{TimingSubsystem}：时序能力的实现，即 \pyapi{exp.timing} 背后的对象。
\end{description}

请注意所有权的清晰划分：\tfocus{时间真相归 \pyapi{sequence}，硬件归 \pyapi{devices}，校准归 \pyapi{\_calibration}，历史归 \pyapi{history}}。读出子系统在工作时会去问设备要原始图、问校准要几何、调 operations 做检测，但它\tfocus{不拥有}这些状态，只是编排它们。

\subsection{一个最小会话的全貌}

\begin{codeblock}[Python]
import Zou_lab_control.neutral_atom as na

session = na.connect(config)

# 会话身上挂着的状态：
session.devices       # DeviceSet：本次连接的硬件
session.sequence      # 当前 PulseSequence（时间真相）
session.history       # list：历次结果的账本

# 会话暴露的实验能力（编排层）：
session.readout       # ReadoutSubsystem
session.timing        # TimingSubsystem
\end{codeblock}

\section{三种后端，同一张会话脸}

本框架最重要的工程价值之一：\tfocus{三种后端共享完全相同的会话表面}。也就是说，无论你连的是离线虚拟设备、本地真实硬件、还是远程 FPGA，你拿到的 \pyapi{NeutralAtomSession} 的接口一模一样。换后端只换 \pyapi{config}，\tfocus{不改一行实验脚本}。

\begin{description}
\item[虚拟后端（离线）] 全部用 \filepath{virtual.py} 里的假设备。没有相机、没有 Vivado、没有网络，照样能 \pyapi{na.connect}、构造序列、跑读出（用合成图）、走完整条流程。这是开发、教学、回归测试的主力。
\item[真实 qCMOS + 本地 RuntimeSequencer] 真实 \pyapi{qcmos} 相机 + 本机的 \pyapi{RuntimeSequencer}。脉冲程序在本地编译并驱动硬件，适合相机和时序硬件都接在同一台机器上的情形。
\item[真实 qCMOS + 远程 FPGA（RemoteSequencer / RPyC）] 真实相机在本地，脉冲流送给\tfocus{远程} FPGA：本地的 \pyapi{RemoteSequencer} 通过 RPyC 调到 \pyapi{sequencer_server.py} 跑的服务（\pyapi{run_server}，jtag-axi 后端），由它经 \pyapi{VivadoAxiStreamerSession} 用 JTAG-to-AXI 操作 FPGA。这是真机实验的典型部署。
\end{description}

\begin{notebox}[换后端只换 config]
因为三种后端共享同一张会话脸，你完全可以\tfocus{先在虚拟后端把整个实验脚本写好测好}，再把 config 切到真机后端原样跑。这正是"同一会话表面"设计带来的最大红利。
\end{notebox}

\begin{codeblock}[Python]
# 同一段实验脚本，三种后端通吃 —— 只有 connect 的 config 不同。
import Zou_lab_control.neutral_atom as na

session = na.connect(config)          # config 决定虚拟 / 本地 / 远程

session.timing.prepare(session.sequence)   # 把时间真相写进（虚拟或真实的）序列器
result = session.readout.run()             # 触发开火 + 读出 + 检测
session.history.append(result)             # 记入账本
\end{codeblock}

\section{走通一次真实实验：完整生命周期}

下面把一次实验从头到尾的生命周期讲清楚。每一步都标注"谁拥有这一步动的状态"，这样你能始终知道自己在和哪一层打交道。

\begin{enumerate}
\item \tfocus{连接}：\pyapi{na.connect(config)} 创建会话，连上设备（\pyapi{session.devices}），初始化空历史、空校准、读出与时序子系统。
\item \tfocus{定义时间}：构造或修改 \pyapi{session.sequence}（一个 \pyapi{PulseSequence}）。这是 timing 层的事，纯数据、无副作用。
\item \tfocus{准备序列}：\pyapi{session.timing.prepare(...)} 把时间真相编译成 \pyapi{RuntimeSequenceProgram} 并写入序列器硬件（本地或远程）。从这一步起才真正"碰硬件"。
\item \tfocus{校准（按需）}：若还没有校准，先做一次得到 \pyapi{TrapCalibration}，让会话持有它。读出要靠它把格点编号映射到相机像素几何。
\item \tfocus{读出}：\pyapi{session.readout} 触发开火、采集\tfocus{原始图}、调用 operations 的检测算法、结合校准产出每个格点的占据判定。
\item \tfocus{记账}：把本次结果追加进 \pyapi{session.history}。
\item \tfocus{分析}：用 \filepath{core/analysis.py}、\filepath{operations/} 对历史结果做离线分析——这些都是纯算法，可随时重跑。
\end{enumerate}

\subsection{一段可运行的最小实验}

\begin{codeblock}[Python]
import Zou_lab_control.neutral_atom as na

# 1) 连接（虚拟后端先跑通，再换真机 config）
session = na.connect(config)

# 2) 定义时间真相：当前序列就是 session.sequence
seq = session.sequence
# ……在此用 PulseSequence 的接口编辑周期 / 通道 / 扫描槽……

# 3) 准备：把时间真相编译并写入序列器
session.timing.prepare(seq)

# 4) 读出：开火 + 采原始图 + 检测 + 结合校准判定占据
result = session.readout.run()

# 5) 记账
session.history.append(result)

# 6) 离线分析历史（纯算法，无副作用，可复现）
#    用 operations/ 与 core/analysis.py 处理 session.history
\end{codeblock}

\begin{notebox}[先在虚拟后端跑这一段]
上面这段脚本在虚拟后端就能完整跑完：序列器是假的、相机是假的、图是合成的，但流程、接口、历史账本都和真机一致。确认逻辑没问题后，把 \pyapi{config} 换成真机后端再跑一次即可。
\end{notebox}

\section{铁律：把边界焊死的不变量}

分层只有靠不变量来强制执行才有意义。本框架有几条\tfocus{绝不能破}的规则。它们看起来是细节，实则是整套架构能换后端、能离线复现、能单测的根本原因。

\begin{description}
\item[capture 只给原始图] \pyapi{capture()} 返回的\tfocus{只有原始图像}，不带任何叠加。格点圆圈、占据标记这些\tfocus{属于读出（readout），不属于相机}。相机驱动的职责到"把像素读回来"为止，再往上一步都不归它。
\item[设备只做硬件动作] devices 层只负责硬件动作，\tfocus{所有算法都住在 operations/ 与 core/}。你不会在相机类里看到拟合，也不会在序列器类里看到检测。
\item[TrapArrayDevice 不知道像素中心] \pyapi{TrapArrayDevice} 暴露 \pyapi{n_sites}，但\tfocus{不}暴露相机像素中心。像素中心来自 \pyapi{TrapCalibration}（在 \filepath{core/calibration.py}）。设备知道"有几个格点"，校准知道"它们在像素上的哪里"——两份知识分属两层。
\end{description}

\begin{notebox}[为什么这几条这么重要]
这三条不变量合起来保证了：相机可以被换成虚拟设备而读出逻辑不变（因为读出不依赖相机去画圈）；算法可以离线对存档图重跑（因为算法不在设备里）；几何校准可以独立保存、加载、复用（因为它不被塞进设备状态）。\tfocus{每当你想在某一层"顺手多做一点别人的事"时，回到这三条铁律——答案几乎总是"不要"。}
\end{notebox}

\section{把分层装进脑子里：一张归属速查表}

最后用一张速查表收尾，帮你在写代码时随时确认"这件事归谁、这份状态在哪"。

\begin{description}
\item[我要改脉冲时长 / 通道延迟 / 扫描] 改 \pyapi{session.sequence}（timing 层，\filepath{sequence.py} / \filepath{pulse_table.py}）。
\item[我要把序列送上硬件] 用 \pyapi{session.timing.prepare(...)}（subsystems 编排 → devices 上传，本地 \pyapi{RuntimeSequencer} 或远程 \pyapi{RemoteSequencer}）。
\item[我要拍一张原始图] 找相机 \pyapi{capture()}（devices/\filepath{qcmos.py}）——只会拿到原始像素，没有圈。
\item[我要知道某格点在图上哪个像素] 问 \pyapi{TrapCalibration}（core/\filepath{calibration.py}），不要去问 \pyapi{TrapArrayDevice}。
\item[我要从图里找原子 / 拟合 / 判定占据] 用 operations/ 的纯算法；在实验流程里由 \pyapi{session.readout} 编排调用。
\item[我跑过的实验结果在哪] \pyapi{session.history}。
\item[我想离线复现一次分析] 拿存档原始图喂给 operations/ 与 core/\filepath{analysis.py}，不需要任何硬件。
\item[我想换后端] 只改 \pyapi{na.connect} 的 \pyapi{config}，脚本一字不改。
\end{description}

记住会话对象是一切的中心：\tfocus{\pyapi{na.connect(config)} 给你一个 \pyapi{NeutralAtomSession}，它拥有 devices / sequence / \_calibration / history，并通过 readout 与 timing 两个子系统把硬件、时间、算法编排成你能跑的实验。}守住分层与不变量，这套系统就能在虚拟、本地、远程三种后端上以同一张脸为你工作。

\chapter{PulseSequence 与 PulseTableState（深入）}

整个时序系统里有两个名字很像、但绝对不能混淆的对象：\tfocus{PulseSequence} 和 \tfocus{PulseTableState}。它们分工不同、拥有的状态不同、生命周期也不同。读完这一章，一个新人应该能说清楚：谁是"物理真相"，谁是"人类编辑模型"，一行边沿表到底是什么意思，以及如何从一个 GUI 表格一路编译到可以上传给 FPGA 的边沿表，并据此跑一次真实实验。

\section{为什么是两个模型}

实验控制有两套截然不同的需求，硬塞进一个类里会两头不讨好，所以代码刻意把它们拆成两层。

\begin{description}
\item[底层真相（PulseSequence）] 定义在 \filepath{Zou\_lab\_control/neutral\_atom/timing/sequence.py}。它是\tfocus{时间正确}的唯一真相：它知道每一个通道在任意时刻的电平。它不关心"周期""人类标签""GUI 可见性"这些概念，只关心一组脉冲（某通道在 start 秒开始、持续 duration 秒为高），以及每通道的整体延迟。它可以把自己编译成一张\tfocus{边沿表}（绝对 tick + 完整输出掩码）。Notebook、PyQt 前端、远程 FPGA 三方共享的就是这一层。

\item[编辑模型（PulseTableState）] 定义在 \filepath{Zou\_lab\_control/neutral\_atom/timing/pulse\_table.py}。它是 GUI 的\tfocus{编译模型}，保存的是给人看的概念：一行一行的\tfocus{周期（PERIOD）}，每个周期有一个时长加一个数字状态向量；还有通道名/标签、可见性、每通道延迟、模拟总线方案（analog-bus plan）、以及扫描槽（scan slots）。它\tfocus{不拥有任何硬件}，它只负责把这些人类概念 \pyapi{compile()} 成一个 PulseSequence / 运行时程序。
\end{description}

\begin{notebox}[一句话区分]
PulseSequence 回答"在第几纳秒，哪些通道是高"；PulseTableState 回答"这张表格里，第几个周期持续多久、哪些灯亮"。前者是物理，后者是界面。PulseTableState 最终\tfocus{下沉}成 PulseSequence，再下沉成边沿表。
\end{notebox}

这样拆分的好处：
\begin{itemize}
\item GUI 可以自由地引入"周期""延迟框""可见通道"这些纯展示概念，而不会污染时间真相。
\item Notebook 用户可以完全绕过表格，直接用 \pyapi{PulseSequence.pulse(...)} 手写脉冲，得到的边沿表和 GUI 编出来的是同一种东西、走同一个上传管线。
\item 只有 PulseSequence 这一层需要为"时间是否落在时钟网格上""掩码怎么打包"负责；PulseTableState 只要保证自己能干净地下沉即可。
\end{itemize}

\section{一行边沿表是什么意思}

边沿表（edge table）是这套系统真正交给硬件的东西。理解它，整条链路就通了。

一行边沿 = \tfocus{"在这个绝对 FPGA tick，把所有输出一次性切换到这个掩码"}。

\begin{itemize}
\item \tfocus{绝对 tick}：从序列起点 t=0 算起的时钟周期计数。50\,MHz 时钟下，一个 tick = 20\,ns；tick=100000 就是 t = 100000 × 20\,ns = 2\,ms。
\item \tfocus{掩码（mask）}：一个整数，按位表示这一刻\tfocus{所有}通道的电平。\tfocus{bit 0 = channels[0]}，bit 1 = channels[1]，依此类推。掩码是"全量快照"，不是"增量"——每一行都重新写满所有通道的当前电平。
\item \tfocus{合并}：如果多个通道在\tfocus{同一时刻}一起变化，它们\tfocus{合并成同一行}的同一个掩码。硬件每个 tick 只换一次输出，所以同一瞬间的所有变化天然合到一行。
\end{itemize}

举一个相机成像预设在 50\,MHz 下编译出来的真实例子。假设通道列表里 ch00=cooling（冷却）、ch03=probe（探测）、ch09=trap（阱）、ch11=emCCD（相机触发位）。编译结果大致是：

\begin{codeblock}[text]
ticks = [0,      100000, 105000, ...]
masks = [513,    512,    2568,   ...]
\end{codeblock}

逐行解读：
\begin{itemize}
\item \pyapi{tick=0, mask=513}。513 的二进制是 1000000001，即 bit0 + bit9 = ch00 + ch09 = \tfocus{cooling + trap} 同时为高。序列一开始就把这两路打开。
\item \pyapi{tick=100000, mask=512}。512 = bit9 = 只剩 ch09=trap。也就是说在 tick 100000（50\,MHz 下 = 2\,ms）那一刻，cooling 关掉了，trap 仍开。注意 cooling 的"关"是\tfocus{隐式}的——掩码里那一位变成 0 就是关。
\item \pyapi{tick=105000, mask=2568}。2568 = bit3 + bit9 + bit11 = ch03 + ch09 + ch11 = \tfocus{probe + trap + emCCD}。这一刻探测光打开、阱保持、相机触发位拉高。emCCD 就是相机触发位——它出现在掩码里，相机就在这一 tick 被触发曝光。
\end{itemize}

\begin{notebox}[为什么掩码是全量而不是增量]
全量掩码让硬件极其简单：每到一个边沿 tick，直接把这个整数怼到输出寄存器，不需要"记住上一刻是什么再做异或"。代价是每行要写满所有位，但通道数固定（62 条），整数装得下，存储也省。
\end{notebox}

\section{从 PulseSequence 到边沿表：compile\_runtime\_program}

把一个 PulseSequence 变成可上传程序的函数是 \pyapi{compile_runtime_program(sequence, *, channels, clock_hz, trigger_channels)}，定义在 \filepath{Zou\_lab\_control/neutral\_atom/devices/sequencer.py}。它返回一个 \tfocus{RuntimeSequenceProgram}。

它内部做的事，按顺序：
\begin{enumerate}
\item 规整通道列表，校验 clock\_hz 为正。
\item 取 \pyapi{sequence.without_repeat()} 得到"一份未展开的基础序列"，再调用 \pyapi{base_sequence.edges(clock_hz=..., channels=...)} 得到 \pyapi{(ticks, masks, channels)}——这就是上面讲的边沿表。
\item 计算 loop\_end\_tick：如果序列设了重复周期 repeat\_period，就用它换算成 tick；否则用最后一条边沿的 tick。这决定了硬件循环回绕的位置。
\item 用 \pyapi{_ensure_final_off_edge(...)} 在循环末尾补一条"全关"边沿，保证回绕前输出干净。
\item 把 sequence、clock\_hz、channels、ticks、masks 等打包，算一个 sha256 截断作为 \pyapi{sequence_id}（用于去重/缓存）。
\item 组装并返回 RuntimeSequenceProgram，里面带上 ticks/masks、duration、trigger\_count（用 trigger\_channels 数触发脉冲个数）、repeat\_forever、loop\_start\_index、loop\_end\_tick、loop\_count 等元数据。
\end{enumerate}

\tfocus{关键点}：边沿表本身\tfocus{不把重复展开}。一个会重复 1000 次的成像循环，边沿表只存一份，配上 loop\_end\_tick 和 loop\_count，让硬件自己回绕。这就是为什么内部先 \pyapi{without_repeat()} 拿基础边沿，再单独带上循环元数据。

一个纯 Notebook、不碰 GUI 的最小例子：

\begin{codeblock}[Python]
from Zou_lab_control.neutral_atom.timing.sequence import PulseSequence
from Zou_lab_control.neutral_atom.devices.sequencer import compile_runtime_program

# 手写脉冲：cooling 从 0 开 2ms；trap 全程开 2.1ms；probe+emCCD 在 2.1ms 触发 0.1ms
seq = (
    PulseSequence(name="imaging")
    .pulse("ch00", start=0.0,    duration=2.0e-3)   # cooling
    .pulse("ch09", start=0.0,    duration=2.1e-3)   # trap
    .pulse("ch03", start=2.1e-3, duration=0.1e-3)   # probe
    .pulse("ch11", start=2.1e-3, duration=0.1e-3)   # emCCD 相机触发
)

prog = compile_runtime_program(
    seq,
    channels=[f"ch{n:02d}" for n in range(62)],   # bit0=ch00, bit1=ch01, ...
    clock_hz=50e6,                                  # 20ns/tick
    trigger_channels=["ch11"],                      # 数相机触发个数
)

print(prog.ticks)          # 绝对 tick 列表
print(prog.masks)          # 与 ticks 对齐的全量掩码列表
print(prog.trigger_count)  # = 1（emCCD 触发了一次）
\end{codeblock}

这里 \pyapi{channels} 的\tfocus{顺序就是位序}：channels[0]→bit0。如果你把通道列表换个顺序，同样的脉冲会编出不同的掩码整数——但物理上等价。

\section{PulseTableState 的状态与 API}

PulseTableState 是 GUI 那张表的内存模型。它拥有的状态（谁拥有什么）：

\begin{description}
\item[periods] 一串周期，每个周期 = 一个时长 + 一个数字状态向量（哪些通道在这个周期为高）。这是表格的每一行。
\item[channels / channel\_labels / visible\_channels] 通道名、给人看的标签、以及哪些通道在 GUI 里可见（纯展示，不影响编译出的物理）。
\item[delays / delay\_units] 每通道的整体延迟，以及它的单位。延迟会在下沉成 PulseSequence 时整体平移该通道。
\item[analog\_buses / analog\_bus\_modes] 模拟总线（DAC 总线）的成员通道，以及每个周期该总线是"edge（写一个值）"还是"hold（保持）"、写什么值。
\item[scan\_slots / scan\_table] 扫描槽（有序，每个绑定到一个 \tfocus{duration} 或 \tfocus{dac} 字段——\pyapi{SCAN_SLOT_KINDS = ("duration", "dac")}，延时不可绑）和扫描表（N\_points × N\_slots）。表达式里用 s0..s\{N-1\} 引用这些槽。
\end{description}

核心方法（都在 \filepath{Zou\_lab\_control/neutral\_atom/timing/pulse\_table.py}）：

\begin{description}
\item[\pyapi{bind_field(kind, target, *, label, unit, nominal)}] 把一个字段绑定到\tfocus{新的扫描槽}，并把该字段原地改写成 \pyapi{s\{i\}}（i 是新槽号），返回槽号。kind ∈ \{duration, dac\}：duration 的 target 是周期下标，dac 的 target 指向某条总线的某个周期。\tfocus{延时不能绑定}——\pyapi{bind_field("delay", ...)} 抛 \pyapi{ValueError("delay is a fixed per-channel value and cannot be scanned")}；延时通过 \pyapi{delays} 设成固定值。\tfocus{幂等}：重复绑定同一个已绑定字段，直接返回它已有的槽号。绑定时会给 scan\_table 每行追加该字段的 nominal 标称值。
\item[\pyapi{unbind_slot(index, *, restore=None)}] 删除第 index 个扫描槽；后面的槽\tfocus{整体降号}（s2→s1…），并重新把每个剩余槽的变量写回对应字段。restore 给被解绑字段一个恢复值。
\item[\pyapi{slot_index_for(kind, target)}] 查某个字段当前绑到了哪个槽，没绑返回 None。bind\_field 的幂等就是靠它实现的。
\item[\pyapi{set_scan_table(rows)} / \pyapi{load_scan_table(path)}] 设置/从文件载入扫描表，会按当前槽数 normalize。表是 N\_points 行 × N\_slots 列。
\item[\pyapi{with_slots_resolved(slots)}] 返回一个\tfocus{非扫描}的副本，把每个槽替换成常量值（时间槽吃 ns、dac 槽吃整数 DAC 码）。用于"每发设一个值"的简洁 Notebook 单点扫描。
\item[\pyapi{analog_bus_plan(name)} / \pyapi{set_analog_bus_mode(...)} / \pyapi{set_bus_value(period, bus, value)}] 读/写某条模拟总线的逐周期方案（edge/hold + 值）。
\item[\pyapi{to_sequence(...)}] 把这张表下沉成一个 PulseSequence（遍历周期，把每段连续为高的区间变成 \pyapi{seq.pulse(...)}，再对每通道套上延迟）。这是 PulseTableState → PulseSequence 的桥。
\item[\pyapi{compile(*, clock_hz, trigger_channels, slots, repeat_forever)}] 把表（解析掉扫描槽后的单点）编译成一个静态运行时程序。内部转去 \pyapi{compile_pulse_table_runtime_program}。
\item[\pyapi{compile_scan(*, clock_hz, scan_table, trigger_channels, repeat_forever)}] 把表连同扫描表一起编译成一个\tfocus{带扫描}的运行时程序（硬件逐点扫）。内部转去 \pyapi{compile_pulse_table_scan_runtime_program}。
\end{description}

\begin{notebox}[槽号是位置而非名字]
扫描槽是\tfocus{有序}的，表达式只能写 s0、s1……不能起别名。所以 \pyapi{unbind_slot} 会触发降号重排——删了 s1，原来的 s2 就变 s0 之后的 s1。这一点在手写表达式或读 scan\_table 列时务必记住：列顺序严格对应 scan\_slots 的顺序。
\end{notebox}

\section{一个完整的 PulseTableState 编译实例}

下面从零搭一张成像表，绑一个扫描槽，分别走\tfocus{单点编译}和\tfocus{扫描编译}两条路，最后说明怎么据此跑实验。

\begin{codeblock}[Python]
from Zou_lab_control.neutral_atom.timing.pulse_table import PulseTableState

# 1) 建表：62 个通道，3 个周期
state = PulseTableState(channels=[f"ch{n:02d}" for n in range(62)])

# 周期 0: cooling+trap 亮 2ms；周期 1: 仅 trap 5us；周期 2: probe+trap+emCCD 亮 0.1ms
# （states 是数字状态向量：哪些通道在该周期为高。具体填法见 GUI / period API）
# 这里假设已经通过界面/构造把三行周期填好，时长分别约 2ms / 5us / 0.1ms。

# 2) 把"周期 0 的时长"绑成一个扫描槽 —— 我们要扫冷却时间
slot = state.bind_field("duration", "0", unit="ns", nominal=2_000_000.0)
print("cooling 时长绑到槽:", slot)   # -> 0，字段现在是 's0'

# 3) 给扫描表：扫 5 个冷却时长（单位 ns），一列对应 s0
state.set_scan_table([[1_000_000.0],
                      [1_500_000.0],
                      [2_000_000.0],
                      [2_500_000.0],
                      [3_000_000.0]])

# 4a) 单点编译：把 s0 解析成一个常量，得到一份静态程序
prog_single = state.compile(
    clock_hz=50e6,
    slots={"s0": 2_000_000.0},   # 这一发用 2ms 冷却
)
print(prog_single.ticks[:4], prog_single.masks[:4])

# 4b) 扫描编译：把整张扫描表交给硬件逐点扫
prog_scan = state.compile_scan(clock_hz=50e6)
print("扫描点数:", len(prog_scan.scan_points))   # -> 5
\end{codeblock}

两条路的区别：
\begin{itemize}
\item \pyapi{compile(...)} 走 \pyapi{slots} 把每个 s\{i\} 钉成常量，得到\tfocus{一份}静态边沿表——适合"我就跑这一个参数"。
\item \pyapi{compile_scan(...)} 把 scan\_table 一起带上，编出的 RuntimeSequenceProgram 里有 scan\_points 和 per-edge 的槽系数，硬件能\tfocus{无缝}地一点一点扫过去，不必每点重新上传。
\end{itemize}

\subsection{编出来之后怎么跑实验}

无论走哪条路，最终拿到的都是一个 RuntimeSequenceProgram，它和 \pyapi{compile_runtime_program} 直接编 PulseSequence 得到的是\tfocus{同一种东西}，因此走同一个上传/运行管线（详见时序与设备相关章节里的 PulseController / 远程 sequencer）。一次最小实验循环大致是：

\begin{enumerate}
\item 在 PulseTableState 里搭好周期、通道、延迟、模拟总线方案；用 \pyapi{bind_field} 绑好你要扫的字段；用 \pyapi{set_scan_table} 给扫描数组。
\item 调 \pyapi{state.compile_scan(clock_hz=50e6)} 得到运行时程序（或单点用 \pyapi{compile}）。
\item 把程序交给设备层（prepare：拉住复位、按行写入、释放；fire：发起脉冲；wait\_done：轮询 done）。
\item 每个扫描点对应一发，相机由掩码里的 emCCD 触发位（如上例 ch11）在对应 tick 自动触发曝光；读出与扫描点一一对应。
\item 跑完回到 safe\_state。
\end{enumerate}

\begin{notebox}[新人最容易踩的三个坑]
\begin{enumerate}
\item \tfocus{别把两个模型搞反}：你在 GUI 里编辑的是 PulseTableState；真正上硬件的是它 compile 出来的边沿表（PulseSequence → RuntimeSequenceProgram）。改了表不 compile，硬件什么都看不到。
\item \tfocus{掩码是位序，不是名字}：513 之所以是 cooling+trap，完全取决于 channels 列表的顺序（bit0=channels[0]）。换了通道顺序，同一物理动作的掩码整数会变。
\item \tfocus{时间必须落在时钟网格上}：50\,MHz 下一切时长/延迟最终都被换算成整数 tick（20\,ns 一格）。不在网格上的时间会在校验阶段报错，而不是被悄悄取整。
\end{enumerate}
\end{notebox}

\chapter{Sequencer 生命周期与远程会话（深入）}

本章把"在笔记本里调一行 \pyapi{pulse.on_pulse()}，原子真的被激光照到"这条完整链路拆开，一层一层讲清楚：谁持有什么状态、每个动作在哪台机器上做什么、纳秒级时序到底由谁掌管。读完本章，一个新人应当能独立启动远程服务、连上它、跑出一次有限次相机采集，并理解为什么这套设计在 Windows + Vivado + JTAG 这种"软件层处处有抖动"的环境里仍然能给出确定的脉冲。

\section{四个动作与两台计算机}

整个控制系统只暴露四个真正的硬件动作：\tfocus{prepare}（准备）、\tfocus{fire}（发射）、\tfocus{wait\_done}（等待完成）、\tfocus{safe\_state}（回到安全空闲）。它们沿着一条固定的链路传递：

\begin{codeblock}[text]
控制计算机 (笔记本/notebook)
  RemoteSequencer            <- 客户端代理, 持有 host/port/channels/clock
      |  RPyC (TCP, 可选 SSL)
      v
FPGA/Vivado 计算机
  SequencerService           <- 有状态服务, 持有 prepared_program + state
      |  callback
      v
  VivadoAxiStreamerSession   <- 常驻 Vivado Tcl/hw_axi 进程, 持有已上传程序
      |  JTAG-to-AXI (axi_bram_ctrl + CTRL mailbox)
      v
  zlc_pulse_streamer_top (FPGA) <- 真正掌管纳秒时序的硬件
\end{codeblock}

\begin{description}
\item[控制计算机] 跑实验脚本/笔记本。它只有"意图"：一份脉冲表（\pyapi{PulseSequence} 或 \pyapi{PulseTableState}）和扫描表。它不直接碰硬件。
\item[FPGA/Vivado 计算机] 接着 FPGA 的那台机器。它跑 RPyC 服务，背后挂一个\tfocus{常驻}的 Vivado Tcl/\pyapi{hw_axi} 进程，通过 JTAG-to-AXI 把 BRAM image 写进 FPGA，并读写 CTRL 寄存器（COMMAND/STATUS mailbox）。
\end{description}

\tfocus{为什么要分两台机器？} 因为 Vivado 启动慢、打开硬件目标（JTAG 链）慢、加载探针文件慢。如果每次"On Pulse"都重新启动一遍，每发一枪都要等几十秒。把这些慢操作集中到\tfocus{服务器启动时}做一次，之后每次 prepare/fire 只是往一个已经活着的 Tcl 进程里喂几行脚本，毫秒级返回。

\section{状态归属：谁持有什么}

理解这套系统的关键，是搞清楚\tfocus{每一层各自缓存了什么状态}，以及这些状态在什么时候失效。

\begin{description}
\item[\pyapi{RemoteSequencer}（客户端）] 持有连接参数（host、port、channels、clock\_hz、trigger\_channels）和 \pyapi{self.\_conn}（RPyC 连接对象）、\pyapi{self.\_last\_program}（上一次 prepare 返回的程序快照）。\pyapi{open()} 在首次使用时建立连接，并用服务器的 \pyapi{snapshot()} 覆盖本地 channels/clock，保证两端通道顺序与时钟一致。
\item[\pyapi{SequencerService}（服务端，进程内状态机）] 持有 \pyapi{self.prepared\_program}（当前已准备好的 \pyapi{RuntimeSequenceProgram}）和 \pyapi{self.state}（\pyapi{idle}/\pyapi{prepared}/\pyapi{running}/\pyapi{done}/\pyapi{timeout}/\pyapi{safe}/\pyapi{aborted}）。它用一把 \pyapi{threading.RLock} 串行化所有动作。\pyapi{fire()} 前会 \pyapi{\_require\_prepared()}，没准备过就直接报错——这是防止"空枪"的护栏。
\item[\pyapi{VivadoAxiStreamerSession}（硬件传输层）] 持有 \pyapi{self.\_process}（常驻 Vivado 子进程）、读取线程、流式回填后台线程，以及当前已载入 FPGA 的程序。它负责把程序打包成 BRAM image、经 JTAG-to-AXI 上传，并驱动 CTRL mailbox。
\item[FPGA（\pyapi{zlc_pulse_streamer_top}）] 持有边沿表（abs ticks + 62 位掩码 + slot 系数，三块并行 BRAM）、repeat 元数据、2-bank 扫描点窗口、bus 段。一旦收到 FIRE 上升沿，\tfocus{枪内的纳秒时序就完全归 FPGA 时钟所有}。
\end{description}

\begin{notebox}[一条铁律：纳秒时序的所有权会"交接"]
prepare 阶段，时序由软件主导（Tcl 把 BRAM image 写进去，多慢都无所谓，因为还没发 FIRE、输出处于安全态）。fire 给出 FIRE 上升沿那一刻，所有权\tfocus{交接}给 FPGA 时钟。此后 Python 的 GIL、RPyC 往返、Windows 调度抖动、Vivado 的延迟、JTAG 的延迟，\tfocus{都不再影响一枪之内的边沿时刻}。这就是为什么"在 Python 里 sleep 然后翻通道"绝对不可接受，而"把整张边沿表交给 FPGA 自己跑"可以做到纳秒确定。
\end{notebox}

\section{prepare：准备但不发射}

\pyapi{prepare(program)} 做四件事，顺序严格：

\begin{enumerate}
\item \tfocus{先发 SAFE}。让 FPGA 回到安全态（输出钳在安全电平），此时往 BRAM 写程序映像不会产生任何脉冲。
\item \tfocus{校验后打包上传}。\pyapi{validate_pulse_streamer_program} 先检查边沿数、扫描点数、tick 位宽、通道数不越界，并校验全局有效 tick 单调（扫描不会把合并后的边沿顺序打乱）；通过后用 \pyapi{host.image.pack_program} 把整张程序打包成 BRAM image（CTRL 标量 + 三块并行边沿表 + 2-bank 扫描窗口 + bus 段），经 JTAG-to-AXI 一次写进 \pyapi{axi_bram_ctrl}。越界或不单调会在写任何 BRAM 之前就报错。
\item \tfocus{arm 扫描 bank}。对流式扫描，把前两个 chunk 装进 ping-pong 的两个 bank，并通过 \pyapi{BANK_READY}/\pyapi{BANK*_CHUNK} 告知引擎哪个 bank 装的是哪一段。
\item \tfocus{发 LOAD}。COMMAND 字先清 0 再置 LOAD，制造干净的上升沿，等回读到 \pyapi{STATUS_LOADED}。准备完成后 FPGA 处于"装好弹、保险还在"的状态。
\end{enumerate}

\tfocus{prepare 绝不启动序列}。它返回后 FPGA 静止，输出维持安全电平。这一点对相机实验至关重要：你可以先 prepare、再武装相机、最后才 fire，三步之间想隔多久都行。

服务端的 \pyapi{SequencerService.prepare} 还做了一层缓存：如果新程序的 \pyapi{sequence_id} 和已准备的完全相同，就跳过 \pyapi{prepare_callback}（即跳过整段上传），直接复用。是否启用由 \pyapi{ZLC_SEQUENCER_CACHE_PREPARED} 控制。

\begin{codeblock}[Python]
from Zou_lab_control.neutral_atom.devices.sequencer import RemoteSequencer

seq = RemoteSequencer(
    host="192.168.1.50",   # FPGA 计算机的可达地址, 不能是 0.0.0.0
    port=18861,
    channels=["cooling", "repump", "imaging", "mot_coil"],
    clock_hz=50e6,          # 20 ns / tick
)
program = seq.prepare(pulse_table_state)   # 上传边沿表, 但不发射
print(program.sequence_id, "triggers:", program.trigger_count)
# 此刻 FPGA 装好弹、保险在, 输出仍是安全电平
\end{codeblock}

\section{fire：一个显式的 start 脉冲，时序所有权交接}

\pyapi{fire()} 发出一个显式的 \tfocus{低-高-低} start 脉冲（\pyapi{0 -> 1 -> 0}）。FPGA 只把\tfocus{同步后的那个上升沿}当作"开始"。用脉冲而不是电平，是为了让"开始"这件事在硬件里有一个干净、唯一、可被时钟同步的事件，不依赖软件维持任何电平。

\pyapi{fire()} 之后，正如前面那条铁律：枪内每条边沿的精确时刻由 FPGA 时钟决定，软件栈的一切抖动都被隔离在枪外。

服务端 \pyapi{fire(sequence_payload)} 还有一个安全检查：如果你传了一个 payload，它会重新编译并比对 \pyapi{sequence_id}，与已 prepare 的不一致就报错——杜绝"准备的是 A，发射的是 B"。

\begin{codeblock}[Python]
seq.fire()                 # 发出 low-high-low start 脉冲, FPGA 开跑
ok = seq.wait_done(timeout=2.0)   # 有限序列: 轮询 STATUS
assert ok, "序列在 2 秒内没有报告 done"
\end{codeblock}

\section{wait\_done：只对有限序列轮询}

\pyapi{wait_done(timeout)} 轮询 \pyapi{STATUS} 寄存器（\pyapi{DONE} 位），等一个\tfocus{有限}序列跑完；对流式扫描它还会跟随 \pyapi{CURSOR} 把空出来的 bank 回填下一段。它有一条硬护栏：

\begin{codeblock}[Python]
# SequencerService.wait_done 里:
if program.repeat_forever and timeout is None:
    raise RuntimeError(
        "sequencer wait_done cannot wait forever for a repeat_forever "
        "program; pass a timeout or stop the pulse.")
\end{codeblock}

也就是说，对一个 \pyapi{repeat_forever=True} 的无限循环程序，\pyapi{wait_done()} 永远等不到 done，必须给 timeout 或者主动 stop。这条护栏在客户端的 \pyapi{PulseController.on_pulse(wait=True)} 里也有对应的、更友好的报错，提示你三种正确组合。

\section{safe\_state：回到安全空闲（Stop Pulse 走这条）}

\pyapi{safe_state()} 把 sequencer 复位到一个安全的空闲态——输出回到安全电平、序列停止。GUI 里的"Stop Pulse"按钮就是调它。在服务端，\pyapi{set_safe_state()} 还会把 \pyapi{prepared_program} 清空、状态标为 \pyapi{safe}，这意味着下一次必须重新 prepare 才能 fire（不能在 safe 之后直接空枪 fire）。

\begin{description}
\item[\pyapi{abort()}] 与 \pyapi{set_safe_state()} 类似，也会作废已准备的程序并调安全回调，区别在语义：abort 是"出错了，紧急停"，safe\_state 是"正常收工，回到待命"。
\end{description}

\section{常驻 Vivado 会话：慢操作只做一次}

\pyapi{VivadoAxiStreamerSession} 是把"慢"集中处理的核心。它在 \pyapi{start()} 时启动\tfocus{一个} Vivado Tcl 子进程（\pyapi{vivado -mode tcl -nolog -nojournal}），然后跑一段公共初始化 Tcl：

\begin{enumerate}
\item 打开硬件目标（\pyapi{open_hw_target}，失败会重试 \pyapi{-jtag_mode on}）。
\item 加载 \filepath{.ltx} 探针文件\tfocus{一次}——这让 Vivado 在已烧录的比特流里找到 \pyapi{jtag_axi} 核（即 \pyapi{hw_axi}）。找不到该核会明确报错，提示当前 FPGA 镜像不是 ZLC pulse-streamer 比特流，或 \filepath{.ltx} 不匹配。
\item 之后整段服务生命周期里\tfocus{复用}这同一个进程、同一个已打开的硬件目标、同一个 \pyapi{hw_axi}。
\end{enumerate}

所以"慢的硬件目标 bring-up 发生在服务器启动时，而不是第一次 On Pulse 时"。每个后续动作（prepare/fire/wait/safe）只是往这个活着的进程喂一段带唯一 marker 的 Tcl（一次 AXI 写或读），读到 marker 即返回。

\begin{notebox}[启动服务器：warm\_start 让你在第一枪之前就付清启动代价]
\pyapi{run_server(backend="jtag-axi", warm_start=True)} 会在接受任何客户端连接之前就 \pyapi{hardware_backend.start()}，把 Vivado 启动 + 打开硬件目标 + 加载探针的几十秒一次性付清。等客户端连上来，第一次 prepare 就已经是热的。
\end{notebox}

\begin{codeblock}[PowerShell]
# 在 FPGA/Vivado 计算机上, 启动常驻服务 (PowerShell)
$env:ZLC_PS_VIVADO_BIN = "C:\Xilinx\Vivado\2025.1\bin\vivado.bat"
$env:ZLC_PS_VIVADO_LTX = "C:\zlc\zlc_pulse_streamer_top.ltx"
python -m Zou_lab_control.neutral_atom.devices.sequencer_server `
    --backend jtag-axi --host 0.0.0.0 --port 18861 `
    --channels cooling repump imaging mot_coil
# 打印 "Starting persistent Vivado JTAG-to-AXI session before accepting clients..."
# 慢启动在这里发生; 之后客户端的 prepare/fire 都是热的
\end{codeblock}

\section{流式回填：扫描点比常驻窗口多时怎么办}

扫描窗口在硬件里是一个 \tfocus{2-bank ping-pong}（\pyapi{BANK_SIZE}=2048 一个 bank，两 bank 共 4096 个常驻扫描点）。当扫描点数 $N$ 超过 4096 时，主机用\tfocus{流式回填}补齐：

\begin{itemize}
\item 引擎依次播放点 \pyapi{0..N-1}，用 \pyapi{(idx/BANK_SIZE) mod 2} 选 bank，并把已消费点数写进 \pyapi{CURSOR}。
\item 主机读 \pyapi{CURSOR}，把\tfocus{游标已经走过的那个 bank}回填下一段 chunk，并更新 \pyapi{BANK_READY}/\pyapi{BANK*_CHUNK}。
\item 引擎只在 \pyapi{BANK_READY} 为真\tfocus{且}该 bank 装的是正确 chunk 时才推进；回填来晚了就\tfocus{停住}（hold，置 \pyapi{STATUS_UNDERFLOW}），\tfocus{绝不}播放错误的点。
\end{itemize}

\pyapi{repeat_forever} 对一个有限（但很大）的流式扫描做反复扫掠：fire 时启动一个后台回填线程，循环把下一段填进空出的 bank；一次扫掠之内是无缝的，扫掠与扫掠之间的接缝是一个短暂的安全 hold。

\begin{notebox}[sequence\_id 缓存让重复 prepare 几乎免费]
一次扫描里若某个点和上一点编译出的程序完全一致，\pyapi{SequencerService.prepare} 通过比对 \pyapi{sequence_id}（程序内容哈希）直接跳过整段上传，第二次 \pyapi{On Pulse} 只是几次控制写。而扫描点本身在硬件里\tfocus{无缝迭代}（仿射 slot 向量 + 2-bank 窗口），并不需要逐点重新 prepare。
\end{notebox}

\section{RemoteSequencer 与 RPyC：客户端代理}

\pyapi{RemoteSequencer} 是控制计算机侧的瘦代理。它本身不碰硬件，只把请求序列化成 JSON 经 RPyC 发给服务端。

\begin{description}
\item[连接] \pyapi{open()} 建立 RPyC 连接（\pyapi{rpyc.connect}，可选 \pyapi{ssl_connect}），并立即调远端 \pyapi{snapshot()}，用服务端的 channels/clock/trigger\_channels 覆盖本地——\tfocus{服务器是通道顺序与时钟的唯一真相来源}。\pyapi{sync_request_timeout=None} 让长动作（如硬件 prepare）不会被客户端超时打断。
\item[host 护栏] 构造时如果 host 是空、\pyapi{0.0.0.0} 或 \pyapi{::} 会直接报错——这些是"监听所有接口"的服务端地址，不是客户端能连的地址。
\item[传输] \pyapi{prepare/fire} 把 \pyapi{PulseSequence}/\pyapi{PulseTableState} 经 \pyapi{timing_payload_to_dict} 转 dict、\pyapi{json.dumps} 后发出；\pyapi{prepare} 收回的程序 JSON 反序列化为 \pyapi{RuntimeSequenceProgram} 存进 \pyapi{\_last_program}。
\end{description}

\begin{codeblock}[Python]
seq = RemoteSequencer(host="192.168.1.50", port=18861,
                      channels=["cooling","repump","imaging","mot_coil"],
                      clock_hz=50e6, connect_on_init=True)
print(seq.snapshot()["remote"]["state"])   # 远端状态机当前态
seq.set_safe_state()                         # 远端回到安全空闲
\end{codeblock}

\section{五个 Sequencer 类各司其职}

仓库里有五个实现 \pyapi{SequencerDevice} 契约的 sequencer，角色互不重叠：

\begin{description}
\item[\pyapi{VirtualSequencer}] 纯软件假人。无硬件，用于离线开发与单元测试，验证编译/契约而不接 FPGA。
\item[\pyapi{RuntimeSequencer}] 本地适配器：直接内嵌一个 \pyapi{SequencerService}，在\tfocus{同一进程}里走完整的 prepare/fire/wait 协议（默认无硬件回调，用 \pyapi{sleep_scale} 模拟时长）。适合在一台机器上端到端联调协议。
\item[\pyapi{ManualSequencer}] 手动/逐步模式，给需要人工介入或单步推进的台架调试用。
\item[\pyapi{RemoteSequencer}] RPyC 客户端代理，连远端真实硬件（本章主角）。
\item[\pyapi{VerilogSequencer}] 走 \filepath{verilog.py} 那条路，把脉冲编译成边沿表/HDL 工件（\pyapi{VerilogBuild}），用于生成/检视而非实时驱动。
\end{description}

\section{repeat\_forever 的两种用法：示波器 vs. 相机}

\pyapi{repeat_forever} 这个旗标的取舍直接决定实验对不对：

\begin{description}
\item[\pyapi{repeat_forever=True}（无限自由跑）] 适合\tfocus{示波器调参}。让 FPGA 不停重复同一序列，你在示波器上实时看波形、转旋钮、立刻看到变化。注意此时 \pyapi{wait_done()} 永远不会返回 done（见前面的护栏），要用 \pyapi{on_pulse(wait=False, repeat_forever=True)}，靠 Stop Pulse 收尾。
\item[\pyapi{repeat_forever=False}（有限）] 适合\tfocus{相机采集}。正确做法：\tfocus{先武装相机}（arm），\tfocus{再生成恰好需要的有限触发序列}，让相机在确定的触发数下读出。
\end{description}

\begin{notebox}[千万不要让 qCMOS 去等一个无限自由跑的循环]
有限相机采集必须用有限序列。如果让 qCMOS 等在一个 \pyapi{repeat_forever=True} 的自由跑循环上，相机不知道何时该停、采了几帧，会拿到错配甚至无限挂起的读出。正确顺序是：arm 相机 $\to$ prepare 有限序列 $\to$ fire $\to$ wait\_done(有限) $\to$ 读图。
\end{notebox}

\begin{codeblock}[Python]
# 一次有限相机采集的完整循环 (控制计算机)
camera.arm(n_frames=expected_triggers)        # 1) 先武装相机
seq.prepare(finite_pulse_with_camera_triggers)# 2) 上传恰好需要的有限触发序列
seq.fire()                                     # 3) start 上升沿, FPGA 接管纳秒时序
ok = seq.wait_done(timeout=expected_duration_s + 0.5)  # 4) 轮询 done
assert ok
frames = camera.read()                          # 5) 读出
seq.set_safe_state()                            # 6) 回到安全空闲
\end{codeblock}

\section{从零跑一次真实实验：完整清单}

把以上各层串起来，一个新人在两台机器上跑出第一枪的步骤：

\begin{enumerate}
\item \tfocus{FPGA 计算机}：跑 \filepath{fpga\textbackslash build\_and\_program.bat}（\filepath{create_project.tcl} 生成工程、编译并烧录 \filepath{zlc_pulse_streamer_top.bit} 与配套 \filepath{.ltx}）。\filepath{.ltx} 必须与比特流来自同一次编译，否则 server 找不到 \pyapi{jtag_axi} 核。
\item \tfocus{FPGA 计算机}：设置 \pyapi{ZLC_PS_VIVADO_BIN} 与 \pyapi{ZLC_PS_VIVADO_LTX}，\pyapi{run_server(backend="jtag-axi", warm_start=True, host="0.0.0.0", port=18861, channels=...)}。等它打印完常驻会话启动信息——慢启动在此付清。
\item \tfocus{控制计算机}：构造 \pyapi{RemoteSequencer(host=<FPGA 机 IP>, port=18861, channels=..., clock_hz=50e6)}，\pyapi{open()} 自动用远端 snapshot 校准通道与时钟。
\item 用 GUI 或脚本做出脉冲表（\pyapi{PulseTableState}），需要扫描就绑 slot、填 scan\_table。
\item 示波器调参阶段：\pyapi{on_pulse(wait=False, repeat_forever=True)} 自由跑，看波形拧旋钮。
\item 正式采集阶段：arm 相机 $\to$ \pyapi{prepare(有限序列)} $\to$ \pyapi{fire()} $\to$ \pyapi{wait_done(timeout)} $\to$ 读图。扫描的后续点靠差分上传，几乎免费。
\item 收工：\pyapi{set_safe_state()}（GUI 的 Stop Pulse），输出回安全电平、作废已准备程序。
\end{enumerate}

\begin{notebox}[一句话记住整章]
慢操作（Vivado/JTAG/探针）在\tfocus{服务器启动}时做一次；prepare 在 reset 拉高、输出钳零的情况下用差分上传写边沿表，绝不发射；fire 给一个 start 上升沿，把纳秒时序的所有权\tfocus{交接}给 FPGA 时钟；wait\_done 只对有限序列有意义；相机采集必须用有限序列、先武装后发射。
\end{notebox}

\chapter{N-slot 扫描模型与硬件无缝扫描（深入）}

本章把“扫描”这件事从头讲透。读完之后，一个新人应该能在 Notebook 里绑定可扫字段、装入扫描表、编译成一个 FPGA 程序，并理解为什么一个一百万点的扫描\tfocus{不会}把边沿表撑大、为什么 DAC 值与周期长度可以\tfocus{在一次连续运行里同时扫描}。（\tfocus{只有 duration 和 dac 两类可扫}；channel 延时不是扫描量，它是一条固定的、有界深度（\pyapi{DELAY_DEPTH}）的逐通道输出延时线，与任意扫描共存——见上一章“固定延时如何与扫描共存”。）

\section{为什么是“命名 slot”而不是 x/y}

旧版本的扫描引擎硬编码了两个扫描轴 \pyapi{x} 和 \pyapi{y}。这套设计已经\tfocus{彻底废除}——现在没有 x，也没有 y。取而代之的是任意多个\tfocus{命名 slot}：\pyapi{s0}、\pyapi{s1}、\pyapi{s2}、……每个 slot 是一个独立的扫描维度。

一个 slot 由三件东西定义：

\begin{description}
\item[kind（种类）] 它绑定的是哪一类字段。只有\tfocus{两种}合法值（\pyapi{SCAN_SLOT_KINDS = ("duration", "dac")}）：\pyapi{"duration"}（某个周期的时长）、\pyapi{"dac"}（某个周期里、某条模拟总线上的 DAC 输出值）。\tfocus{channel 延时不在其中}——它是每条 channel 的固定值，\pyapi{bind_field("delay", ...)} 会报错。
\item[target（目标）] 在该种类里具体指哪一个字段。对 \pyapi{duration} 是周期下标（字符串化的整数，如 \pyapi{"1"}）；对 \pyapi{dac} 是 \pyapi{"<bus>@<period_index>"} 形式，例如 \pyapi{"da_dipole@2"}（第 2 个周期、\pyapi{da\_dipole} 总线上的值）。
\item[unit / nominal（单位与标称值）] 时间类 slot 带单位（\pyapi{ns}/\pyapi{us}/\pyapi{ms}/\pyapi{s}）；\pyapi{dac} slot 的“单位”是\tfocus{原始 10 位 DAC 码}，没有电压换算。\pyapi{nominal} 是绑定那一刻字段的当前值，作为默认/回填值。
\end{description}

slot 的编号就是它的 s 变量：第 \pyapi{i} 个 slot 在所有表达式里就是字符串 \pyapi{s\{i\}}（\pyapi{slot\_var(0) == "s0"}）。这一点很重要——\tfocus{slot 的序号即变量名}，所以删除一个 slot 会让后面的 slot 全部重新编号（见 \pyapi{unbind\_slot}）。

\section{绑定字段：bind\_field 与 scan dot}

把一个普通字段变成“可扫描的符号字段”，叫做\tfocus{绑定}。绑定后，该字段在内部不再是一个数字，而是被改写成符号 \pyapi{s0}/\pyapi{s1}/……，参与表达式求值。

两种绑定途径，效果完全一致：

\begin{itemize}
\item \tfocus{GUI}：点击字段右侧那个嵌入的小圆点（scan dot）。圆点把字段染成淡橙底、深橙字，并标上 slot 编号（橘色）；字段随即被禁用为符号。再点一次解除。
\item \tfocus{API}：调用 \pyapi{state.bind_field(kind, target)}，返回新 slot 的整数下标。
\end{itemize}

状态的所有权很清晰：\filepath{Zou\_lab\_control/neutral\_atom/timing/pulse\_table.py} 里的 \pyapi{PulseTableState} 持有 \pyapi{scan\_slots}（有序的 \pyapi{ScanSlot} 列表）和 \pyapi{scan\_table}（数据）。GUI 只是这个状态的视图；编译器只读这个状态。

\pyapi{bind\_field} 的关键语义：

\begin{itemize}
\item \tfocus{幂等}：重复绑定同一个 \pyapi{(kind, target)} 返回已有下标，不会新建 slot（所以 GUI 里重复点 dot 不会清零或重排编号——编号保留）。
\item 绑定时若不给 \pyapi{nominal}，会从字段当前值自动读取（\pyapi{\_read\_field\_nominal}）。
\item 绑定后会把该字段改写成符号：\pyapi{duration} → 周期时长变成 \pyapi{"s0"}；\pyapi{dac} → 该周期该总线段变成 \pyapi{\{"mode": "edge", "value": "s0"\}}。（\pyapi{delay} 不能绑定——它是固定值，\pyapi{bind_field("delay", ...)} 抛 \pyapi{ValueError}。）
\item 已有的扫描表每一行会自动追加一列（用 \pyapi{nominal} 填充），保持 \pyapi{scan\_table} 列数 = slot 数。
\end{itemize}

\begin{codeblock}[Python]
from Zou_lab_control.neutral_atom.timing.pulse_table import PulseTableState

state = PulseTableState.demo()  # 任意已搭好的脉冲表

# 1) 扫第 1 个周期的时长（按 s0）
s0 = state.bind_field("duration", "1", unit="us")        # 返回 0

# 2) 扫第 2 个周期里 da_dipole 总线上的 DAC 值（按 s1）
s1 = state.bind_field("dac", "da_dipole@2")              # 返回 1

# cooling 的延时是一个固定的每通道值，不进扫描表（也不能绑定）
state.delays["cooling"] = 500.0                          # 固定 500 ns 输出延时
# state.bind_field("delay", "cooling")  # -> ValueError: delay is a fixed per-channel value

print(state.scan_slots)      # 两个 ScanSlot，kind 分别是 duration/dac
\end{codeblock}

\begin{notebox}[s 变量只能引用 slot]
时间表达式里只能出现 \pyapi{s0..s\{N-1\}}，没有 x/y 也没有别的自由符号。一个时长可以写成依赖多个 slot 的\tfocus{仿射}表达式（如 \pyapi{"s0 + 2*s1"}），但必须是 slot 的一次线性组合——见后文“限制”。
\end{notebox}

\subsection{解绑与重新编号}

\pyapi{state.unbind\_slot(index, restore=...)} 删除一个 slot。因为\tfocus{编号即变量名}，删除中间 slot 会把它后面的所有 slot 往前移：原来的 \pyapi{s2} 变成 \pyapi{s1}，依此类推，所有引用它们的字段会被自动重写到新编号（\pyapi{\_renumber\_slot\_bindings}）。\pyapi{restore} 给出解绑后字段恢复成的字面值（不给则时长回到 1000 ns、DAC 回到 hold）。

\section{scan\_table：N\_points × N\_slots 的数据矩阵}

扫描的“数据”和扫描的“结构”是分开的。结构是 \pyapi{scan\_slots}（哪些字段被扫）；数据是 \pyapi{scan\_table}：一个 \tfocus{N\_points 行 × N\_slots 列} 的二维数组。

\begin{itemize}
\item \tfocus{行} = 一个扫描点（一次相机曝光对应一行）。
\item \tfocus{列 j} = slot \pyapi{s\{j\}} 在该点的取值，\tfocus{用该 slot 的显示单位}：时间类 slot 用 ns（或绑定时的单位，内部按 \pyapi{UNITS\_TO\_NS} 折算成 ns），\pyapi{dac} 类用\tfocus{带符号 DAC 值}（-512..+511，0 = 0 V；线上自动 +512 变 offset-binary 码）。
\end{itemize}

装入数据有几种方式：

\begin{description}
\item[从文件] \pyapi{state.load\_scan\_table(path)} 支持 \pyapi{.npy} / \pyapi{.csv} / \pyapi{.txt}。文件里就是一个 N×N\_slots 的矩阵，列顺序必须和 slot 顺序一致。
\item[直接给数组] \pyapi{state.set\_scan\_table(rows)}，会按当前 slot 数归一化（补齐/截断列）。
\item[GUI Scan 页] 在 Scan 标签页里用代码生成参数（每个 slot 一段 Python），或点 Load 从文件载入。
\end{description}

\begin{codeblock}[Python]
import numpy as np

# 两列对应 s0(us 周期时长)、s1(DAC 码)
table = np.array([
    [10.0,    0],     # 点 0：周期=10us，DAC=0
    [20.0,  256],     # 点 1
    [30.0,  768],     # 点 2
    [40.0, 511],      # 点 3：DAC 正满量程
])
state.set_scan_table(table)

# 或从文件
# state.load_scan_table("scan_params.npy")
\end{codeblock}

\begin{notebox}[单位约定再强调一次]
时间列写你看到的数（ns/us/ms/s 取决于该 slot 的 unit）；DAC 列写\tfocus{带符号 10 位值 -512..+511}（0 = 0 V，+511 正满量程、-512 负满量程）。电压到 LSB 的换算由你在写表前自己完成；线上 offset-binary 码由编译器自动 +512 生成。
\end{notebox}

\section{主机如何编译：一个仿射模板 + 一张流式点表}

这是整套设计的核心，也是它省 LUT 的根本原因。主机\tfocus{不}为每个扫描点各生成一份边沿表。它只生成两样东西：

\begin{enumerate}
\item \tfocus{一份仿射边沿模板}：所有边沿的“基准 tick”加上每个边沿对每个 slot 的系数。
\item \tfocus{一张流式扫描点表}：就是 \pyapi{scan\_table} 转成整数码（\pyapi{scan\_points}）。
\end{enumerate}

运行时，FPGA 对每个扫描点用同一份模板，按下式实时算出每条边沿的有效 tick：

\begin{codeblock}[text]
effective_tick = base_tick + (Σ_j coeff_j * slot_j) >> COEFF_FRAC_BITS
\end{codeblock}

其中 \pyapi{COEFF\_FRAC\_BITS = 8}（即 \pyapi{scan\_coeff\_frac\_bits}）。系数用 8 位定点小数表示，所以斜率可以是非整数。\pyapi{slot\_j} 是当前扫描点里 slot j 的整数值（时间类已折算成 tick 量纲，dac 类是 DAC 码）。

由此得到本设计最重要的性质：

\begin{notebox}[大扫描不会撑大边沿表]
边沿表的长度只由\tfocus{模板复杂度}（有多少条边沿、多少周期/通道）决定，\tfocus{与扫描点数无关}。扫 10 个点和扫 100 万个点，边沿模板一模一样；多出来的只是流式点表里多几行整数。存储是 O(edges) + O(points × slots)。
\end{notebox}

编译入口是 \pyapi{state.compile\_scan(clock\_hz=...)}（薄封装），它转调 \filepath{Zou\_lab\_control/neutral\_atom/devices/sequencer.py} 里的 \pyapi{compile\_pulse\_table\_scan\_runtime\_program}，返回一个 \pyapi{RuntimeSequenceProgram}。该程序对象携带：

\begin{description}
\item[slot\_count / slot\_kinds] slot 数量与每个 slot 的 kind。
\item[tick\_slot\_coeffs] 每条边沿对每个 slot 的系数（行=边沿，列=slot）。
\item[loop\_end\_slot\_coeffs] 循环结束 tick 对各 slot 的系数（让循环边界也随扫描平移）。
\item[scan\_points] 整数化的扫描点表（list[list[int]]）。
\item[scan\_coeff\_frac\_bits] 即 8。
\item[ticks / masks] 基准边沿模板（绝对 tick + 62 位通道掩码）。
\item[bus\_segments] 模拟总线段（含 DAC 无缝扫描所需的 value\_select 与段 tick 系数）。
\end{description}

\begin{codeblock}[Python]
program = state.compile_scan(clock_hz=50e6)   # 50 MHz → 20 ns/tick

print(program.slot_count)            # 2
print(program.slot_kinds)            # ['duration', 'dac']
print(program.scan_coeff_frac_bits)  # 8
print(len(program.scan_points))      # 4（点数）
print(program.tick_slot_coeffs)      # 每条边沿 × 2 个系数
print(program.channel_delays)        # 固定延时作为常数携带（不在 slot 里）
\end{codeblock}

FPGA 拿到这个程序后，\tfocus{无缝地}遍历所有扫描点：每点对应模板的一次实例化，每点发\tfocus{一次相机触发}，点与点之间不停机。这就是“硬件无缝扫描”。

\section{DAC 值 + 时长：同时扫两者（并叠加一个固定延时）}

这是相对早期版本最大的突破。早期 DAC 值的扫描和周期时长的扫描是互斥的（怕段的 tick 不跟着时长平移而失步）。现在不互斥了：\tfocus{DAC 值与周期时长可以在同一次无缝运行里同时扫描}。再叠加一个\tfocus{固定}的通道延时也完全没问题——延时不是扫描量，它作为常数携带，作为一条延时线加在输出上、\tfocus{与两类扫描均匀组合}。

实现的两个支柱：

\begin{enumerate}
\item \tfocus{段 tick 也仿射发射}：每个模拟总线段的起止 tick 不再是固定值，而是 \pyapi{base + 系数·slot}，用和边沿同一套 \pyapi{effective\_tick} 公式。于是当你扫一个周期的时长时，落在该周期后面的 DAC 段会\tfocus{自动跟着平移}，不会失步。
\item \tfocus{value\_select 选值来源}：每个总线段带一个选择位。\pyapi{0} 表示用字面值；\pyapi{j+1} 表示运行时从 scan slot \pyapi{j} 读取低位作为 DAC 码。绑定 \pyapi{kind="dac"} 时，编译器把对应段的 value\_select 指向那个 slot。
\end{enumerate}

绑定 DAC 用 \pyapi{target = "<bus>@<period_index>"}，扫描表对应列写\tfocus{带符号 10 位 DAC 值（-512..+511，0 = 0 V）}。下面是一个\tfocus{同时扫 duration 与 dac}、并给一条 channel 设\tfocus{固定}延时的可运行端到端例子：

\begin{codeblock}[Python]
import numpy as np
from Zou_lab_control.neutral_atom.timing.pulse_table import PulseTableState

state = PulseTableState.demo()

# --- 同时绑定两个可扫维度 ---
s0 = state.bind_field("duration", "1", unit="us")   # s0: 第1周期时长(us)
s1 = state.bind_field("dac", "da_dipole@2")         # s1: 周期2 上 da_dipole 的 DAC 值

# --- 一个固定的（不可扫）每通道延时，叠加在输出上 ---
state.delays["cooling"] = 150.0                     # 固定 150 ns 延时线（上限 DELAY_DEPTH≈±40us）

# --- 扫描表：4 个点，列顺序 = s0(us), s1(DAC码) ---
state.set_scan_table(np.array([
    [10.0,    0],
    [20.0,  256],
    [30.0,  768],
    [40.0, 511],
]))

# --- 一次编译，得到一个无缝扫描程序 ---
program = state.compile_scan(clock_hz=50e6)

assert program.slot_count == 2
assert program.slot_kinds == ["duration", "dac"]
assert len(program.scan_points) == 4
# DAC 段带 value_select 指向 s1；该段 tick 随 s0 平移（仿射系数非零）
# cooling 的固定延时作为常数落在 program.channel_delays（不在任何 slot 里）
\end{codeblock}

运行这个程序时，FPGA 对四个点依次实例化模板：每点周期 1 的时长、周期 2 上 \pyapi{da\_dipole} 的输出码都按该行取值，DAC 段的位置随时长平移；cooling 的固定延时在每一点都把它的输出整体推后 150\,ns（一条延时线，深度上限 \pyapi{DELAY_DEPTH}），每点打一次相机触发——全程不停机。

\begin{notebox}[谁拥有什么状态：生命周期一览]
\begin{itemize}
\item \pyapi{PulseTableState}（pulse\_table.py）拥有 \pyapi{scan\_slots} + \pyapi{scan\_table} + 脉冲结构。GUI 是它的视图。
\item \pyapi{compile\_scan / compile\_pulse\_table\_scan\_runtime\_program}（sequencer.py）把状态\tfocus{只读}编译成 \pyapi{RuntimeSequenceProgram}（仿射模板 + 整数点表 + 总线段）。
\item RTL（\filepath{fpga/pulse\_streamer/zlc\_edge\_streamer.v}）按 \pyapi{effective\_tick} 公式实时实例化模板，无缝遍历 \pyapi{scan\_points}，每点一相机触发；固定延时由引擎作为逐通道输出延时线施加。
\end{itemize}
\end{notebox}

\section{限制与编译器拒绝的情形}

仿射模板 + 流式点表很强，但有它换来省 LUT 的代价。编译器会主动拒绝违反前提的输入——这是好事，它在主机端就把会失步的程序挡下来：

\begin{description}
\item[时间表达式必须仿射] 所有 duration 表达式必须是 slot 的一次线性组合（\pyapi{s0 + 2*s1} 可以，\pyapi{s0*s1} 或 \pyapi{s0**2} 不行）。非仿射的时序会被拒绝。（延时不参与扫描表达式——它是固定值。）
\item[边沿顺序不能在扫描中翻转] 如果某个扫描点会让两条边沿的先后顺序交换（模板的边沿排序在不同点之间不一致），编译器拒绝。换句话说，整段扫描里边沿的相对顺序必须稳定。
\item[ramp 端点不能是被扫的值] DAC 无缝扫描里，唯一的限制是：\tfocus{斜坡（ramp）的端点不能是被扫描的那个值}。需要扫描该值时，把该段改成\tfocus{边沿 hold}（\pyapi{mode="edge"}）而不是 ramp。
\item[DAC 列是码不是电压] 再说一遍：\pyapi{dac} slot 的扫描表列是带符号 10 位值 -512..+511（0 = 0 V），线上才是 offset-binary 码。
\end{description}

\begin{notebox}[单点 Notebook 扫描的便捷写法]
如果你只想一发设一个值（不走硬件无缝遍历），用 \pyapi{state.with\_slots\_resolved(\{"s0": 20.0, "s2": 768\}\)}：它返回一个\tfocus{非扫描}副本，把每个 slot 替换成常量（时间类给 ns，dac 给整数码），\pyapi{scan\_slots} 和 \pyapi{scan\_table} 清空。适合在循环里逐点手动设值的简易扫描。
\end{notebox}

\chapter{真实硬件 Runbook（深入）}

本章是一份"从零到完成一次真实成像"的完整操作手册。它面向第一次接触本系统的人：把整套控制链路拆成几层，逐层说明\tfocus{谁拥有什么状态}、什么时候做什么、命令长什么样、出问题时往哪里看。读完本章并照做，你应当能在两台真实计算机上把 FPGA 编译、烧写、起服务、从 notebook 连上去打一次相机成像脉冲。

\section{系统全景：两台计算机、四个层次}

整套系统跑在\tfocus{两台物理计算机}上，中间用 RPyC 网络连接：

\begin{description}
\item[控制/qCMOS 计算机] 跑实验 notebook 和相机。这台机器上 \pyapi{na.connect(...)} 创建一个 \tfocus{NeutralAtomSession} 会话对象；它管理相机、trap array、sitemap、以及一个\tfocus{远程} sequencer 句柄。这台机器\tfocus{不需要}装 Vivado。
\item[FPGA/Vivado 计算机] 物理上连着 Artix-7（xc7a35tfgg484-2）开发板和 Vivado。这台机器跑两件事：一次性的\tfocus{编译+烧写}（\filepath{fpga\textbackslash build\_and\_program.bat}），以及常驻的\tfocus{脉冲流服务器}（\filepath{fpga\textbackslash run\_server.bat}）。服务器内部持有一个常驻 Vivado \pyapi{hw_axi} 会话，通过 JTAG-to-AXI 把脉冲程序映像写进 FPGA 的 BRAM。
\end{description}

从上到下，请把控制链路理解成四层，每层有明确的\tfocus{状态所有者}：

\begin{enumerate}
\item \tfocus{Notebook / Session 层}（控制机）：\pyapi{exp = na.connect(...)} 返回的会话对象。它拥有"当前序列"（\pyapi{exp.sequence}）、相机配置、标定结果。
\item \tfocus{RemoteSequencer / RPyC 层}：会话里的 sequencer 句柄是一个网络代理。它把 \pyapi{prepare/fire/wait\_done/safe\_state} 这些动作通过 RPyC 发到 FPGA 机的 \tfocus{SequencerService}。
\item \tfocus{Server / Vivado 会话层}（FPGA 机）：\filepath{run\_server.bat} 起的进程持有一个常驻 Vivado \pyapi{hw_axi} 会话（jtag-axi 后端），是真正会经 JTAG-to-AXI 写 BRAM、改 FPGA 内存的人。它拥有"已经烧进去的 bit 文件"和"当前已载入的程序"。
\item \tfocus{硬件层}：FPGA 里跑的 \filepath{zlc\_pulse\_streamer\_top}（边沿表 + 1-tick 预取 + 2-bank 流式扫描 + 仿射扫描引擎 + 模拟总线），以及 62 路物理输出。它拥有"正在跑的脉冲"和真实的电平。
\end{enumerate}

\begin{notebox}[一句话记忆]
Notebook 说"我要打这个序列"，Session 把它编译成边沿表/扫描表通过 RPyC 发给 FPGA 机的 Vivado 会话，Vivado 会话把它打包成 BRAM image 经 JTAG-to-AXI 写进 FPGA，FPGA 自己按时钟把电平打出来。相机的外部触发也是这 62 路里的一根线。
\end{notebox}

\begin{notebox}[路径长度：一定要短]
把整个仓库放在\tfocus{很短的路径}下，例如 \filepath{D:\textbackslash ZLC}。原因是 Vivado 2019 的 debug-core 会在工程目录里生成很深的临时路径，而 Windows 对路径长度敏感；仓库放得太深会让综合/实现在写 debug-core 临时文件时莫名失败。这不是建议，是硬约束——默认工程目录是 \filepath{fpga\textbackslash build\textbackslash pulse\_streamer}，再加上仓库本身的深度，很容易超界。
\end{notebox}

\section{第一步：在两台机器上安装}

两台机器都要装。安装脚本是 \filepath{install\_requirements.bat}，它做四件事：把 \pyapi{requirements.txt} 装进\tfocus{当前选中的解释器}、用 \pyapi{pip install -e .} 以可编辑方式装上本仓库、注册一个 Jupyter kernel、并把这个解释器的完整路径\tfocus{记到} \filepath{.zlc\_python\_path} 文件里。

\begin{codeblock}[PowerShell]
# 在仓库根目录（例如 D:\ZLC）打开 PowerShell，两台机器都要做一次
cd D:\ZLC

# 方式 A：让脚本自动从 VSCode/Jupyter kernel 推断解释器
.\install_requirements.bat

# 方式 B：直接指定你要用的 python.exe（最稳）
.\install_requirements.bat C:\Users\you\miniconda3\envs\zlc\python.exe
\end{codeblock}

记住 \filepath{.zlc\_python\_path} 这个文件的作用：之后 \filepath{pulse\_gui.bat}（脉冲编辑器 GUI）和 \filepath{fpga\textbackslash run\_server.bat}（服务器）都会\tfocus{优先读取同一个解释器}，这样保证 GUI、服务器、notebook 用的是\tfocus{同一套依赖}。这是避免"我这边能跑、那边报 ImportError"的关键设计。

\begin{description}
\item[控制/qCMOS 机] 装完后，在 VSCode 里选 kernel "Python (Zou lab control)"，重启 Jupyter kernel，就可以跑 notebook。
\item[FPGA/Vivado 机] 装完后还要确认 Vivado 能被找到。脚本会在 \filepath{C:\textbackslash Xilinx\textbackslash Vivado\textbackslash <版本>\textbackslash bin\textbackslash vivado.bat} 和 \filepath{D:\textbackslash Xilinx\textbackslash ...} 下自动探测多个版本。如果你的 Vivado 不在默认路径，显式指定：
\end{description}

\begin{codeblock}[PowerShell]
# 仅当自动探测找不到 Vivado 时才需要设置（指向 vivado.bat）
$env:ZLC_PS_VIVADO_BIN = "C:\Xilinx\Vivado\2019.2\bin\vivado.bat"
\end{codeblock}

\section{第二步（FPGA 机）：编译 + 烧写}

所有 FPGA 侧编译/烧写都走一个入口：\filepath{fpga\textbackslash build\_and\_program.bat}。它有几个模式：

\begin{description}
\item[\texttt{-{}-check}] \tfocus{离线自检}：不连硬件、不用真实引脚，只做 HDL 综合 + 容量自查。第一次部署、或改过 RTL/Tcl 之后，先跑这个。它核对 top 是 62 路 wrapper，BRAM image 几何（边沿/扫描/总线区基址）和 \pyapi{host.image.solve_capacity} 在 35T 上解出的 4096 边沿 + 4096 常驻扫描点 + 流式一致。
\item[（无参数）] \tfocus{编译并烧写}：跑 create-project Tcl 生成工程、综合、实现、生成 bit/ltx，然后烧进 FPGA。
\item[\texttt{-{}-build-only}] 只编译，不烧。
\item[\texttt{-{}-program-only}] 用已有的 bit/ltx 直接烧。
\item[\texttt{-{}-diagnose}] 列出 Vivado 看得到的硬件 target/device，确认 JTAG 线缆连上了。
\end{description}

\begin{codeblock}[PowerShell]
cd D:\ZLC

# 1) 先离线自检：HDL 综合 + 容量
.\fpga\build_and_program.bat --check

# 2) 自检过了，连上板子，确认能看到硬件
.\fpga\build_and_program.bat --diagnose

# 3) 完整编译 + 烧写（耗时较长，是真综合+实现）
.\fpga\build_and_program.bat
\end{codeblock}

默认工程落在 \filepath{fpga\textbackslash build\textbackslash pulse\_streamer}；编译产物 \tfocus{bit 文件和 ltx 探针文件}在它的 \filepath{pulse\_streamer.runs\textbackslash impl\_1} 子目录下，文件名是 \filepath{zlc\_pulse\_streamer\_top.bit} 和 \filepath{zlc\_pulse\_streamer\_top.ltx}。

\begin{notebox}[XDC 引脚映射与 .ltx 探针文件]
默认 XDC 是 address-switch 板的引脚映射（\filepath{fpga\textbackslash board\_config\textbackslash board.xdc}，见 \filepath{fpga/board\_config/README.md}）。换板子/换线缆映射时用 \pyapi{$env:ZLC_PS_XDC} 覆盖。create-project Tcl 会生成 \pyapi{jtag_axi} + \pyapi{axi_bram_ctrl} + BRAM 这套数据通路。这一点很关键：服务器是\tfocus{靠 JTAG-to-AXI 写 BRAM 控制 FPGA}的，所以服务器加载的 \filepath{.ltx} 必须和烧进去的 bit 是\tfocus{同一次编译}产生的，否则找不到 \pyapi{jtag_axi} 核。
\end{notebox}

\section{相机成像的物理输出（必须记住的四根线）}

成像序列只用到四路物理输出，请把这四条对应关系刻在脑子里——它们既是 XDC 里的引脚，也是 notebook 里 \pyapi{configure_imaging} 会驱动的通道：

\begin{description}
\item[trap（光阱）] \filepath{ch09} $\rightarrow$ 引脚 \tfocus{M17}
\item[cooling（冷却）] \filepath{ch00} $\rightarrow$ 引脚 \tfocus{F15}
\item[probe（探测）] \filepath{ch03} $\rightarrow$ 引脚 \tfocus{N15}
\item[emCCD / 相机外触发] \filepath{ch11} $\rightarrow$ 引脚 \tfocus{M13}（这根线就是 qCMOS 的外部触发）
\end{description}

Session 在编译成像序列时会自动做这个映射：它检查 sequencer 报上来的通道里是否同时含有 \filepath{ch00}/\filepath{ch03}/\filepath{ch09} 且有触发通道，若是就把 trap/cooling/probe/trigger 绑到 ch09/ch00/ch03/触发通道。所以你在 notebook 里写 \pyapi{configure_imaging()}，底层就是在这四根线上排脉冲。

\begin{notebox}[时钟陷阱：50 MHz / 20 ns]
默认时钟是 \tfocus{50 MHz}，最小步进 \tfocus{20 ns}（一个 tick）。如果你在示波器上看到脉冲长度\tfocus{正好是设定值的两倍}，几乎可以肯定是某一层还在按\tfocus{旧的 100 MHz}假设算 tick。逐层检查 control（notebook 传的 \pyapi{clock-hz}）、server（\filepath{run\_server.bat} 的 \pyapi{ZLC_PS_CLOCK_HZ}）、GUI（编辑器里的时钟设置）三处是否都是 50 MHz。三处必须一致。
\end{notebox}

\section{第三步（FPGA 机）：启动脉冲流服务器}

烧好 bit 之后，在 FPGA 机上常驻一个服务器进程，它就是 RPyC 的服务端。入口是 \filepath{fpga\textbackslash run\_server.bat}。

先\tfocus{干跑一次配置检查}（不真正起服务，只打印它将要用的全部配置），确认 bit/ltx/XDC/通道/触发/时钟/容量都对：

\begin{codeblock}[PowerShell]
cd D:\ZLC
.\fpga\run_server.bat --check-config
\end{codeblock}

它会打印类似这样的清单（请逐行核对）：

\begin{codeblock}[text]
Host:    0.0.0.0:18861
Backend: jtag-axi
Project: ...\fpga\build\ps
Bit:     ...\impl_1\zlc_pulse_streamer_top.bit
LTX:     ...\impl_1\zlc_pulse_streamer_top.ltx
Channels: ch00 ch01 ... ch61
Clock:   50000000 Hz
Capacity: max_edges=4096 bank_size=2048 (4096 resident) + UNBOUNDED streaming
\end{codeblock}

确认无误后，正式起服务：

\begin{codeblock}[PowerShell]
.\fpga\run_server.bat
\end{codeblock}

默认它在 \tfocus{host 0.0.0.0、port 18861} 上监听，后端是 \tfocus{jtag-axi}（持久 Vivado \pyapi{hw_axi} 会话）。窗口会一直挂着——这是预期行为，\tfocus{不要关掉它}。常用环境变量覆盖：

\begin{description}
\item[\texttt{ZLC\_PS\_HOST} / \texttt{ZLC\_PS\_PORT}] 监听地址/端口，默认 \pyapi{0.0.0.0} / \pyapi{18861}。
\item[\texttt{ZLC\_PS\_CLOCK\_HZ}] 时钟，默认 \pyapi{50000000}。
\item[\texttt{ZLC\_PS\_XDC}] 覆盖引脚映射 XDC（通道名/触发通道是从 XDC 推断的）。
\item[\texttt{ZLC\_PS\_VIVADO\_BIN}] Vivado 不在默认路径时指定 \filepath{vivado.bat}。
\item[\texttt{ZLC\_FPGA\_SERVER\_PYTHON}] 强制服务器用某个 python.exe（否则读 \filepath{.zlc\_python\_path}）。
\end{description}

\begin{notebox}[server 启动失败最常见的原因]
\filepath{run\_server.bat} 在起服务前会硬性要求找到 \tfocus{.ltx 探针文件}（server 靠它在比特流里定位 \pyapi{jtag_axi} 核）。如果报 "no Vivado .ltx probe file was found" 或 ".ltx ... does not exist"，说明你还没烧 bit、或者 bit 和 ltx 不在它找的默认 impl\_1 路径下。解决：先跑 \filepath{build\_and\_program.bat}，或用 \pyapi{$env:ZLC_PS_VIVADO_LTX} 指到 Vivado Program Device 时实际用的那个 Probes 文件。
\end{notebox}

\section{第四步（控制机）：从 notebook 连上去}

服务器跑起来后，回到\tfocus{控制/qCMOS 机}的 notebook。连接用 \pyapi{na.connect}，第一个位置参数是\tfocus{配置名}（真实硬件用 \pyapi{"remote_template"}），用 \pyapi{sequencer={...}} 传远程地址，\pyapi{open_devices=True} 让它立刻打开相机等真实设备。

\begin{codeblock}[Python]
import Zou_lab_control.neutral_atom as na

exp = na.connect(
    "remote_template",
    sequencer={"host": "192.168.0.42", "port": 18861},  # FPGA 机的地址/端口
    open_devices=True,
)
\end{codeblock}

\pyapi{exp} 就是 Session 对象，它拥有当前实验的所有状态。接下来配置成像序列、做飞行前检查、再真正打：

\begin{codeblock}[Python]
# 1) 配置一次成像序列。默认 trigger_width=20us, pre_trigger=100us。
#    exposure 不传则用相机当前曝光；load=True 表示序列里包含装载阶段。
seq = exp.timing.configure_imaging(exposure=20e-3)

# 2) 飞行前检查：编译 verilog、做安全/容量校验，并在失败时直接抛异常。
exp.timing.preflight().raise_if_failed()

# 3) 看一眼将要打的序列（可选，画在 notebook 里）
exp.timing.plot_sequence()
\end{codeblock}

\pyapi{configure_imaging} 的参数都是关键字参数：

\begin{description}
\item[\texttt{exposure}] 相机曝光（秒）。传了就同时把相机 \pyapi{configure(exposure=...)}，不传则沿用相机当前值。
\item[\texttt{load}] 是否在序列里包含原子装载阶段，默认 \pyapi{True}。
\item[\texttt{trigger\_width}] 给相机的外触发脉冲宽度，默认 \pyapi{20e-6}（20 us）。这就是 emCCD/ch11/M13 那根线上的脉冲。
\item[\texttt{pre\_trigger}] 触发前的提前量，默认 \pyapi{100e-6}（100 us）。
\end{description}

\pyapi{preflight()} 返回一个 \tfocus{PreflightReport}：它带 \pyapi{ok}、\pyapi{errors}、\pyapi{warnings}、序列表、设备快照、以及编译出的 verilog 信息。\pyapi{raise_if_failed()} 在 \pyapi{ok} 为假时把所有 error 拼成异常抛出——\tfocus{务必}在真正打硬件前调它。想看细节可以 \pyapi{print(exp.timing.preflight().summary())}。

完成检查后，就可以真正采集 / 读出（具体的 capture/readout API 见设备手册与主手册对应章节，这里强调的是它们\tfocus{此时}才会触发真实硬件动作）。一次最小的"连接到打出脉冲"流程到此就闭环了。

\begin{notebox}[一次完整闭环的状态流转]
\pyapi{configure_imaging} 在 Session 里生成 \pyapi{exp.sequence}（控制机内存）→ \pyapi{preflight} 把它编译成边沿表/扫描表（仍在控制机）→ 采集时通过 RemoteSequencer 经 RPyC 把表发到 FPGA 机的 Vivado 会话 → Vivado 会话 \pyapi{prepare}（保持复位、用 prog\_we 翻转逐行写入、释放）→ \pyapi{fire}（启动脉冲）→ \pyapi{wait\_done}（轮询 done）→ \pyapi{safe\_state}。相机在 emCCD 那根线上被同一份序列触发，图像回到控制机。
\end{notebox}

\section{排错手册（按现象查）}

\begin{description}

\item[现象：示波器上脉冲长度正好是设定值的两倍] 时钟假设错了。某层还按 100 MHz 算 tick，而硬件是 50 MHz（20 ns/tick）。逐一核对：notebook 连接/序列的时钟、\filepath{run\_server.bat} 的 \pyapi{ZLC_PS_CLOCK_HZ}（应为 \pyapi{50000000}）、GUI 编辑器里的时钟。三处必须都是 50 MHz。

\item[现象：综合/实现在写 debug-core 临时文件时失败，路径报错] 仓库放得太深。把整个仓库移到很短的路径（如 \filepath{D:\textbackslash ZLC}）重试。Vivado 2019 的 debug-core 临时路径对长度敏感。

\item[现象：\texttt{run\_server.bat} 报找不到 .ltx 探针文件] 还没烧 bit，或 bit/ltx 不在默认 \filepath{impl\_1} 下。先 \filepath{build\_and\_program.bat}（或 \texttt{-{}-program-only}），或用 \pyapi{$env:ZLC_PS_VIVADO_LTX} 指向实际烧写时用的 Probes 文件。务必保证 ltx 和 bit 来自\tfocus{同一次编译}。

\item[现象：找不到 Vivado] 自动探测覆盖 2019.1–2025.2 的默认 \filepath{C:}/\filepath{D:} 路径。不在其中就设 \pyapi{$env:ZLC_PS_VIVADO_BIN} 指向 \filepath{vivado.bat}。

\item[现象：notebook 报 ImportError，但 GUI/服务器都正常] 三处用的解释器不是同一个。重跑 \filepath{install\_requirements.bat} 把仓库装进 notebook 选中的那个 kernel，并确认 \filepath{.zlc\_python\_path} 指向同一解释器。\filepath{pulse\_gui.bat} 和 \filepath{run\_server.bat} 都会优先读这个文件。

\item[现象：HDL 综合 / 容量自检失败] RTL、top、image 几何漂移了。\filepath{build\_and\_program.bat -{}-check} 会报哪个 BRAM 区段或宽度对不上 \pyapi{host.image.solve_capacity} 解出的契约（边沿 4096、扫描 bank 2048、mask=62、tick=32）。\pyapi{host.image} 是单一真相源——RTL localparam 和 create-project Tcl 都从它派生，改回契约值后重新 create-project。

\item[现象：连不上 FPGA 机 / RPyC 超时] 确认 \filepath{run\_server.bat} 窗口还开着且在监听 \pyapi{0.0.0.0:18861}；notebook 的 \pyapi{sequencer={"host": ..., "port": 18861}} 里 host 是 FPGA 机的真实 IP；防火墙放行 18861。先在 FPGA 机上 \texttt{-{}-check-config} 确认它确实起在期望端口。

\item[现象：相机不触发，但其它脉冲正常] 检查 emCCD 那根线：通道 \filepath{ch11} $\rightarrow$ 引脚 \tfocus{M13}，它就是 qCMOS 外触发。确认 XDC 里这根线没被改、\pyapi{configure_imaging} 的 \pyapi{trigger_width}（默认 20 us）足够让相机识别。服务器配置检查里 Channels/触发通道是从 XDC 推断的，若触发通道推断失败 \filepath{run\_server.bat} 会直接报错退出。

\item[现象：JTAG 看不到板子] 跑 \filepath{build\_and\_program.bat -{}-diagnose} 列出 Vivado 能看到的 hw target/device。看不到就检查 JTAG 线缆、板子供电、以及是否有别的进程占着 hw\_server。

\end{description}

\begin{notebox}[日常启动顺序（背下来）]
（FPGA 机）确认 bit 已烧 → \filepath{run\_server.bat -{}-check-config} 核对 → \filepath{run\_server.bat} 起服务并\tfocus{保持窗口开着}；（控制机）\pyapi{na.connect("remote_template", sequencer=..., open_devices=True)} → \pyapi{configure_imaging(...)} → \pyapi{preflight().raise_if_failed()} → 采集。只有改过 RTL/引脚映射时才需要重新 \filepath{build\_and\_program.bat}。
\end{notebox}

\chapter{Pulse API：脚本控制与延时校准}

GUI 之外，同一套接口可以在 notebook / 脚本里直接驱动 sequencer——\tfocus{GUI 与 API 走完全相同的路径}（同一个 \pyapi{sequencer.prepare/fire/set_safe_state}，同一个 payload 记录），所以两边永远可以互相同步。核心对象是 \pyapi{PulseController}（\pyapi{bind_pulse(sequencer, pulse)} 创建）：

\begin{codeblock}[Pulse API 一览]
from Zou_lab_control.neutral_atom.devices.sequencer import RemoteSequencer, bind_pulse
from Zou_lab_control.neutral_atom import PulseTableState

seq = RemoteSequencer(host="...", port=18861, channels=[f"ch{i:02d}" for i in range(62)])
pulse = bind_pulse(seq, PulseTableState.load("pulses/T.json"))   # 或 pulse.load_pulse(path)

pulse.on_pulse()                      # 编译 + 上传 + FIRE（=GUI 的 On Pulse）
pulse.off_pulse()                     # 停止并回安全态（=GUI 的 Stop Pulse）
pulse.set_channel_delay("ch11", 250)  # 设 emCCD 延时 250 ns（链式可连写）
pulse.get_channel_delay("ch11")       # 读回（ns）
pulse.set_scan_table([[10],[20]])     # 上硬件扫描表
pulse.save_pulse("pulses/T2.json")    # 存盘（GUI 可直接 Load）
state = pulse.synced_state()          # 读回"设备上实际应用的"PulseTableState
\end{codeblock}

\section{延时校准完整示例}

逐通道延时（事件调度，上限约 $\pm 42.9$ s，毫秒级 emCCD 延时直接可用）的典型校准循环——扫一组延时、每个点 on/off 一次、采一组计数，最后用前端接口画图：

\begin{codeblock}[channel delay calibration]
import numpy as np
import Zou_lab_control.frontend as zf

delays_ns = np.linspace(-400, 400, 41)
signal = []
for d in delays_ns:
    pulse.set_channel_delay("ch11", float(d)).on_pulse()
    signal.append(measure())        # 你的采集函数（相机计数/PMT 等）
    pulse.off_pulse()

fig = zf.plot(delays_ns, np.asarray(signal), kind="1d",
              labels=("delay (ns)", "counts", ""), display=False)
fig.fig.savefig("delay_calibration.png")
best = delays_ns[int(np.argmax(signal))]
pulse.set_channel_delay("ch11", float(best)).on_pulse()   # 应用最优延时
\end{codeblock}

\begin{notebox}
脚本改完设备后回到 GUI，点 \tfocus{Sync} 即把"设备上实际应用的脉冲"拉回编辑器（服务端在每次成功 prepare 时记录源 payload，GUI 与 \pyapi{pulse.synced_state()} 读的是同一份）。反过来，GUI 在 RUNNING 状态下被编辑会让 On Pulse 按钮带上星号（\pyapi{On Pulse*}）、标题栏状态点变橙（UNSYNCED）——再点一次即应用。
\end{notebox}

\chapter{修改指南}

\section{我想改默认时钟}
真实硬件默认 50 MHz 是因为当软件/server 假设 100 MHz 时测到的脉冲长了一倍。若换到经过验证的 100 MHz 时钟源，需要同时改 XDC 时钟约束和 \pyapi{ZLC_PS_CLOCK_HZ=100000000}。所有 duration/delay/scan 值都按 \pyapi{1e9 / clock_hz} 对齐到 tick 网格。

\section{我想加一条新 channel 或新预设}
预设是 \filepath{pulses/} 下的 \pyapi{PulseTableState} JSON，是前端/API 数据，不是硬件工程。新预设可以用 GUI 编辑保存，也可以用 API 构造后 \pyapi{state.save(path)}。channel 顺序由 server 从 XDC 推断；GUI 隐藏只是视图操作，上传时缺失的 channel 补 off。

\section{我想加一个扫描参数}
不要把扫描网格展开成很多 GUI 列或很多 prepared 表。绑定一个新字段（\pyapi{bind_field}），给 \pyapi{scan_table} 增加一列，硬件会处理剩下的事。\pyapi{dac} slot 现在也走硬件无缝扫描（见下一节“扫描 DAC 值”），但不能和\tfocus{被扫的持续时间}同一次上传（扫 duration 会移动周期起点，让模拟段的绝对 tick 失步——和\tfocus{延时}扫描则可以同用）。需要扫描模拟 \pyapi{ramp} 的端点时，按扫描点逐个 prepare/fire。

\section{我想改文档}
本手册的正文模板在 \filepath{Zou_lab_control/neutral_atom/content/manual_templates/main_manual_zh.texbody}。改完正文后用 \pyapi{build_main_manual()} 重新生成 PDF（见 \filepath{docs/MAINTAINER_NOTES.md} 第 7 节）。

\chapter{常见问题}

\section{脉冲长度是设定的两倍}
确认 server、GUI、notebook 都显示 50 MHz / 20 ns。仍按 100 MHz 编译时，设定 1 ms 会在 50 MHz 硬件上输出 2 ms。

\section{On Pulse 没报错但示波器无信号}
按顺序检查：FPGA 已 program 当前 bit/LTX；\pyapi{--check-config} 指向同一套 bit/LTX/XDC；GUI 打开的是正确预设；示波器接的是 \pyapi{ch09/ch00/ch03/ch11} 对应的物理输出；脉冲是有限 shot 还是 repeat forever，scope 触发是否设对。

\section{GUI 只显示少数 channel}
这是正常的。GUI 可见列只为降低编辑噪声；上传宽度由 server channel list 决定，缺失 channel 自动补 off。

\section{扫描不通过}
确认每个扫描点的值都是 20 ns 的整数倍；确认时间表达式是 slot 的仿射函数；确认边沿顺序在所有扫描点下不交换。否则拆 chunk 或逐点 prepare。


\end{document}
